% \iffalse meta-comment
%
% Copyright (C) 2024
% 
% -----------------------------------
%
% This work may be distributed and/or modified under the
% conditions of the LaTeX Project Public License, either
% version 1.3c of this license or (at your option) any later
% version. This version of this license is in
%    http://www.latex-project.org/lppl/lppl-1-3c.txt
% and the latest version of this license is in
%    http://www.latex-project.org/lppl.txt
% and version 1.3 or later is part of all distributions of
% LaTeX version 2020/11/27 or later.
%
% \fi
%
% \iffalse
%<cls>\RequirePackage{expl3,l3keys2e}
\ProvidesExplClass{suepthesis}
%<cls>{2024-10-07}{3.8.0}{SUEP Thesis Templates}
%
%<*driver>
\ProvidesFile{suepthesis.dtx}[2024/10/07 0.0.1 SUEP Thesis Templates]
\documentclass[letterpaper]{l3doc}
\usepackage{dtx-style}

\DisableCrossrefs

\RecordChanges
\begin{document}
  \DocInput{\jobname-doc.tex}
  \PrintChanges
  \def\indexname{代码索引}
  %%%%% \PrintIndex
\end{document}
%</driver>
% \fi
%
% \begin{macro}
% \section{实现细节}
%
%    \begin{macrocode}
%<*package>
%    \end{macrocode}
%
% Identify the internal prefix (\LaTeX3 \pkg{DocStrip} convention).
%    \begin{macrocode}
%<@@=suepthesis>
%    \end{macrocode}
%
%    \begin{macrocode}
%</package>
%    \end{macrocode}
%
% \subsection{suepthesis.cls 文档类}
%
%
% \subsubsection{全局变量与临时变量}
%
% 定义全局变量。
% \begin{variable}{\g_@@_thesis_type_int}
% 论文类型，取值从 1 开始，分别对应:
%  \begin{enumerate}
%      \item 本科生毕业设计（论文）
%      \item 硕士学位论文
%      \item 博士学位论文
%  \end{enumerate}
%    \begin{macrocode}
\int_new:N \g_@@_thesis_type_int
%    \end{macrocode}
% \end{variable}
%
% \begin{variable}{\g_@@_thesis_degree_type_int}
% 学位类型，取值从 1 开始，分别对应:
%  \begin{enumerate}
%      \item 学术
%      \item 专业
%  \end{enumerate}
%    \begin{macrocode}
\int_new:N \g_@@_thesis_degree_type_int
%    \end{macrocode}
% \end{variable}
%
% \begin{variable}{\g_@@_head_zihao_int}
% 页眉字号。研究生论文使用 5 号字，本科生论文使用 4 号字。
%    \begin{macrocode}
\int_new:N \g_@@_head_zihao_int
%    \end{macrocode}
% \end{variable}
%
% \begin{variable}{\g_@@_twoside_bool}
% 是否双面打印。默认单面打印。
%    \begin{macrocode}
\bool_new:N \g_@@_twoside_bool
%    \end{macrocode}
% \end{variable}
%
% \begin{variable}{\g_@@_thesis_type_english_bool}
% 是否为英文模板。目前只有本科生全英文专业的模板会将此变量设置为 true。
%    \begin{macrocode}
\bool_new:N \g_@@_thesis_type_english_bool
%    \end{macrocode}
% \end{variable}
%
% \begin{variable}{\g_@@_blind_mode_bool}
% 是否为盲审模式。默认为 false。
%    \begin{macrocode}
\bool_new:N \g_@@_blind_mode_bool
%    \end{macrocode}
% \end{variable}
%
% \begin{variable}{\g_@@_quirks_mode_bool}
% 是否兼容更符合北理工官方模板或规范，但不太符合一般排版要求的模式。默认为 false。
% 目前此选项包括：
% \begin{itemize}
%   \item \pkg{biblatex} 的专利格式不再使用国标格式，而采用北理工自定义格式。
% \end{itemize}
%    \begin{macrocode}
\bool_new:N \g_@@_quirks_mode_bool
%    \end{macrocode}
% \end{variable}
%
% \begin{variable}{\g_@@_label_divide_char_tl}
% 用于分隔标签的字符。默认为「-」或者「.」。
%    \begin{macrocode}
\tl_new:N \g_@@_label_divide_char_tl
%    \end{macrocode}
% \end{variable}
%
% \begin{variable}{\l_@@_right_seq, \l_@@_left_seq}
% 定义临时变量。
%    \begin{macrocode}
\seq_new:N \l_@@_right_seq
\seq_new:N \l_@@_left_seq
%    \end{macrocode}
% \end{variable}
%
% \subsubsection{辅助函数与常量}
%
% \begin{macro}[added=2023-05-06,updated=2023-07-04]{\@@_secret_info:nn,\@@_secret_info:N,\@@_secret_info:n}
% 普通模式下显示参数一，盲审模式下显示参数二。
% \begin{macrocode}
\cs_new:Npn \@@_hide:n #1 {
  \g_@@_const_substitute_symbol_tl
}

\cs_new:Npn \@@_secret_info:nn #1 #2 {
  \bool_if:nTF \g_@@_blind_mode_bool {
    #2
  } {
    #1
  }
}

\cs_new:Npn \@@_secret_info:N #1 {
  \@@_secret_info:nn {#1} {\tl_map_function:NN #1 \@@_hide:n }
}
\cs_new:Npn \@@_secret_info:n #1 {
  \@@_secret_info:nn {#1} {\tl_map_function:NN {#1} \@@_hide:n }
}
% \end{macrocode}
%
% \end{macro}
%
% \begin{macro}[added=2023-03-16]{\@@_get_const:N}
% 获取标题、章节、表格、图形等的常量名称。
% 会区别英文模式和中文模式。
%   \begin{macrocode}
\cs_new:Npn \@@_get_const:N #1 {
  \@@_if_thesis_english:TF {
    \use:c {c_@@_label_ #1 _en_tl}
  } {
    \use:c {c_@@_label_ #1 _tl}
  }
}
%   \end{macrocode}
% \end{macro}
%
% \begin{macro}[added=2023-03-16]{\@@_set_english_mode:}
% 设置为英文模式。
%   \begin{macrocode}
\cs_new:Npn \@@_set_english_mode: {
  \bool_gset_true:N \g_@@_thesis_type_english_bool
}
%   \end{macrocode}
% \end{macro}
%
% \begin{macro}{\tl_if_empty:xTF,\seq_set_split:Nnx}
% 生成变体。
%    \begin{macrocode}
\cs_generate_variant:Nn \tl_if_empty:nTF {x}
\cs_generate_variant:Nn \seq_set_split:Nnn {Nnx}
%    \end{macrocode}
% \end{macro}
%
% \begin{macro}{\@@_same_page:}
% 取消换页。
%    \begin{macrocode}
\cs_new:Npn \@@_same_page: {
  \let\clearpage\relax
  \let\cleardoublepage\relax
}
%    \end{macrocode}
% \end{macro}
% \begin{macro}{\@@_if_graduate:TF}
% 是否为研究生学位论文。
%    \begin{macrocode}
\cs_new:Npn \@@_if_graduate:TF #1#2 {
    \int_compare:nNnTF {1} < {\g_@@_thesis_type_int}
      {#1}
      {#2}
  }
%    \end{macrocode}
% \end{macro}

% \begin{macro}{\@@_if_thesis_int_type:nT,\@@_if_thesis_int_type:nTF}
% 是否某一特定模板。
%    \begin{macrocode}
\cs_new:Npn \@@_if_thesis_int_type:nTF #1#2#3 {\int_compare:nNnTF {\g_@@_thesis_type_int} = {#1} {#2} {#3}}
\cs_new:Npn \@@_if_thesis_int_type:nT #1#2 {\@@_if_thesis_int_type:nTF {#1} {#2} {}}
\cs_new:Npn \@@_if_thesis_int_type:nF #1#2 {\@@_if_thesis_int_type:nTF {#1} {} {#2}}
%    \end{macrocode}
% \end{macro}

% \begin{macro}{\@@_if_thesis_english:T,\@@_if_thesis_english:TF}
% 是否为英文模板，这里包括全英文专业和研究生模板的英文模式。
%    \begin{macrocode}
\cs_new:Npn \@@_if_thesis_english:TF #1#2 {\bool_if:nTF {\g_@@_thesis_type_english_bool} {#1} {#2}}
\cs_new:Npn \@@_if_thesis_english:T #1 {\@@_if_thesis_english:TF {#1}{}}
%    \end{macrocode}
% \end{macro}
%
% \begin{macro}{\@@_if_bachelor_thesis:TF,\@@_if_bachelor_thesis:T,\@@_if_master_thesis:TF,\@@_if_doctor_thesis:TF}
% 是否为本科、硕士、博士学位论文。
%    \begin{macrocode}
\cs_new:Npn \@@_if_bachelor_thesis:TF #1#2 {\int_compare:nNnTF {\g_@@_thesis_type_int} = {1} {#1} {#2}}
\cs_new:Npn \@@_if_bachelor_thesis:T #1 {\@@_if_bachelor_thesis:TF {#1} {}}
\cs_new:Npn \@@_if_master_thesis:TF #1#2 {\int_compare:nNnTF {\g_@@_thesis_type_int} = {2} {#1} {#2}}
\cs_new:Npn \@@_if_doctor_thesis:TF #1#2 {\int_compare:nNnTF {\g_@@_thesis_type_int} = {3} {#1} {#2}}
%    \end{macrocode}
% \end{macro}
%
% \begin{macro}{\@@_if_academic_thesis:TF}
% 是否为学术论文。
%    \begin{macrocode}
\cs_new:Npn \@@_if_academic_thesis:TF #1#2 {\int_compare:nNnTF {\g_@@_thesis_degree_type_int} = {1} {#1} {#2}}
%    \end{macrocode}
% \end{macro}
%
% \begin{variable}{\c_@@_thesis_type_clist}
% 定义论文类型的列表。
%    \begin{macrocode}
\clist_const:Nn \c_@@_thesis_type_clist
    { bachelor, master, doctor }
%    \end{macrocode}
% \end{variable}
%
% \begin{variable}{\c_@@_thesis_degree_type_clist}
% 定义学位类型的列表。
%    \begin{macrocode}
\clist_const:Nn \c_@@_thesis_degree_type_clist
    { academic, professional }
%    \end{macrocode}
% \end{variable}
%
% \begin{variable}{\c_@@_publication_modes_clist}
% 定义「攻读学位期间发表论文与研究成果清单」管理方式。
%    \begin{macrocode}
\clist_const:Nn \c_@@_publication_modes_clist
    { biblatex, custom }
%    \end{macrocode}
% \end{variable}
%
% \begin{macro}{\@@_define_label:nn,\@@_define_label_by_thesis_type:nnn,\@@_define_label:nnn,\@@_define_label_by_thesis_type:nnnn}
% 定义常量（标签）的辅助函数。
%    \begin{macrocode}
\cs_new_protected:Npn \@@_define_label:nn #1#2
  { \tl_const:cn { c_@@_label_ #1 _tl } {#2} }

\cs_new_protected:Npn \@@_define_label_by_thesis_type:nnn #1#2#3
  {
    \tl_const:cn { c_@@_ #1 _label_ #2 _tl } {#3}
  }

\cs_new_protected:Npn \@@_define_label:nnn #1#2#3
  {
    \tl_const:cn { c_@@_label_ #1    _tl } {#2}
    \tl_const:cn { c_@@_label_ #1 _en_tl } {#3}
  }

\cs_new_protected:Npn \@@_define_label_by_thesis_type:nnnn #1#2#3#4
  {
    \tl_const:cn { c_@@_ #1 _label_ #2 _tl } {#3}
    \tl_const:cn { c_@@_ #1 _label_ #2 _en_tl } {#4}
  }
%    \end{macrocode}
% \end{macro}
%
% \begin{macro}{\smallgap:}
% 标签文字之间的间距。
%    \begin{macrocode}
\cs_new:Npn \smallgap: {
  \hspace{0.45ex}
}
%    \end{macrocode}
% \end{macro}
%
% \begin{macro}{\boxempty:}
% 空的选框。
%    \begin{macrocode}
\cs_new:Npn \boxempty:
{
  \makebox[1em][l]
  {
    % 为保证与打了勾的一致，也需套盒子
    \makebox[0pt][l]
    {
      % 默认比基线略高，向下降降
      \raisebox{-1pt}{$\square$}
    }
  }
}
%    \end{macrocode}
% \end{macro}
%
% \begin{macro}{\boxcheck:}
% 打了勾的选框。
%    \begin{macrocode}
\cs_new:Npn \boxcheck:
{
  \makebox[1em][l]
  {
    \makebox[0pt][l]
    {
      % 默认比基线略高，向下降降
      \raisebox{-1pt}{$\square$}
    }
    $\checkmark$
  }
}
%    \end{macrocode}
% \end{macro}
%
% \begin{macro}{\label_space:}
% 标签与内容之间的空白间距。
%    \begin{macrocode}
\cs_new:Npn \label_space: {
  \@@_if_bachelor_thesis:T {
    \quad
  }
}
%    \end{macrocode}
% \end{macro}
%
% \begin{variable}{\c_@@_label_code_tl,\c_@@_label_udc_tl,
%   \c_@@_label_classification_tl,\c_@@_label_classified_level_tl,\c_@@_label_type_tl}
% 没有对应英文的常量。
%    \begin{macrocode}
\clist_map_inline:nn
  {
    {code} {代码},
    {udc} {UDC分类号：},
    {classification} {中图分类号：},
    {classified_level} {密级},
    {type} {种类},
    {special_type} {特别类型},
    {cross_research} {交叉研究方向},
    {international_student_ugp} {政府项目留学生},
  }
  {\@@_define_label:nn #1}
%    \end{macrocode}
% \end{variable}
%
% \begin{variable}{\c_@@_bachelor_label_xxx_tl}
% 本科毕设的常量。
%    \begin{macrocode}
\clist_map_inline:nn
  {
    {originality} {上海电力大学学位论文原创性声明},
    {originality_clause} {本人郑重声明：所呈交的学位论文，是本人在导师的指导下，
    独立进行研究工作所取得的成果。 除文中已经注明引用的内容外， 
    本论文不包含任何其他个人或集体已经发表或撰写过的作品成果。 
    对本文的研究做出重要贡献的个人和集体， 均已在文中以明确方式标明。
    本人完全意识到本声明的法律结果由本人承担。},
    {authorization} {上海电力大学学位论文版权使用授权书},
    {authorization_clause} {本人完全了解北京理工大学有关保管、使用毕业设计（论文）的规定，
    其中包括：\circled{1}~学校有权保管、并向有关部门送交本毕业设计（论文）的原件与复印件；
    \circled{2}~学校可以采用影印、缩印或其它复制手段复制并保存本毕业设计（论文）；
    \circled{3}~学校可允许本毕业设计（论文）被查阅或借阅；\circled{4}~学校可以学术交流为目的，
    复制赠送和交换本毕业设计（论文）；\circled{5}~学校可以公布本毕业设计（论文）的全部或部分内容。},
    {originality_author_signature}
      {本人签名：\hspace{40mm}日\hspace{2.5mm}期：\hspace{13mm}年\hspace{8mm}月\hspace{8mm}日},
    {originality_supervisor_signature}
      {指导老师签名：\hspace{40mm}日\hspace{2.5mm}期：\hspace{13mm}年\hspace{8mm}月\hspace{8mm}日},
  } {\@@_define_label_by_thesis_type:nnn {bachelor} #1}
%    \end{macrocode}
% \end{variable}
% \begin{variable}{\c_@@_graduate_label_xxx_tl}
% 研究生模板的常量。
%    \begin{macrocode}
\clist_map_inline:nn
  {
    {originality} {上海电力大学学位论文原创性声明},
    {originality_clause} {本人郑重声明：所呈交的学位论文，是本人在导师的指导下，
    独立进行研究工作所取得的成果。 除文中已经注明引用的内容外， 
    本论文不包含任何其他个人或集体已经发表或撰写过的作品成果。 
    对本文的研究做出重要贡献的个人和集体， 均已在文中以明确方式标明。
    本人完全意识到本声明的法律结果由本人承担。},
    {authorization} {上海电力大学学位论文版权使用授权书},
    {authorization_clause} {本学位论文作者完全了解学校有关保留、 使用学位论文的规定， 
    同意学校保留并向国家有关部门或机构送交论文的复印件和电子版， 允许论文被查阅和借阅。
    本人授权上海电力大学可以将本学位论文的全部或部分内容编入有关数据库进行检索， 
    可以采用影印、缩印或扫描等复制手段保存和汇编本学位论文。\\
    },
    {originality_author_signature}
      {签\qquad 名：\hspace{40mm}日\hspace{2.5mm}期：\hspace{30mm}\quad},
    {originality_supervisor_signature}
      {导师签名：\hspace{40mm}日\hspace{2.5mm}期：\hspace{30mm}\quad},
  } {\@@_define_label_by_thesis_type:nnn {graduate} #1}
%    \end{macrocode}
% \end{variable}
%
% \begin{variable}{\c_@@_graduate_label_xxx_tl,\c_@@_graduate_label_xxx_en_tl}
% 研究生模板的中英常量。
%    \begin{macrocode}
\clist_map_inline:nn
  {
    {author} {学\quad 位\quad 申\quad 请 \quad 人} {Candidate~Name},
    {supervisor} {指\quad 导\quad 教\quad 师} {Faculty~Mentor},
    {major} {学\quad 科 \quad 专 \quad 业} {Major},
    {degreeclass} {学\quad 位·\quad 类\quad 别} {Degree~Applied},
    {institute}{所\quad 属\quad 院\quad 系}{Degree~by},
    {chairman} {答辩委员会主席} {Chair,~Thesis~Committee},
    {finalizationdate} {论文定稿日期} {Time}
  } {\@@_define_label_by_thesis_type:nnnn {graduate} #1}
%    \end{macrocode}
% \end{variable}
%
% \begin{variable}{\c_@@_label_xxx_tl,\c_@@_label_xxx_en_tl}
% 常用的中英常量。
%    \begin{macrocode}
\clist_map_inline:nn
  {
    {school} {学\qquad 院} {School},
    {major} {专\qquad 业} {Degree},
    {course} {课程名称} {Course},
    {class} {班\qquad{}级} {Class},
    {author} {学生姓名} {Author},
    {student_id} {学\qquad 号} {Student~ID},
    {supervisor} {指导教师} {Supervisor},
    {co_supervisor} {校外指导教师} {Co-Supervisor},
    {teacher} {任课教师} {Teacher},
    {semester} {上课学期} {Semester},
    {keywords} {关键词：} {Key~Words:~},
    {toc} {目\label_space: 录} {Table~of~Contents},
    {abstract} {摘\label_space: 要} {Abstract},
    {conclusion} {结\label_space: 论} {Conclusions},
    % 附录部分的总标题
    {appendix} {附\label_space: 录} {Appendices},
    {ack} {致\label_space: 谢} {Acknowledgements},
    {figure} {插\label_space: 图} {Illustrations},
    {table} {表\label_space: 格} {Tables},
    % 附录下各部分编号的前缀
    {appendix_prefix} {附录} {Appendix},
    {reference} {参考文献} {References},
    {university} {上海电力大学} {Beijing~Institute~of~Technology},
    {publications} {攻读学位期间发表论文与研究成果清单}
      {Publications~During~Studies},
    % TODO: Not so sure about the translation.
    {resume} {作者简介} {Author~Biography},
    {symbols} {主要符号对照表} {Nomenclature},
    {algo} {算法} {Algorithm},
    {them} {定理} {Theorem},
    {lem} {引理} {Lemma},
    {prop} {命题} {Proposition},
    {cor} {推论} {Corollary},
    {axi} {公理} {Axiom},
    {defn} {定义} {Definition},
    {conj} {猜想} {Conjecture},
    {exmp} {例} {Example},
    {case} {情形} {Case},
    {rem} {注} {Remark},
    {fig} {图} {Figure},
    {tab} {表} {Table},
    {equ} {式} {Equation},
  }
  {\@@_define_label:nnn #1}
%    \end{macrocode}
% \end{variable}
%
% \begin{variable}{\c_@@_bachelor_thesis_header_clist,\c_@@_bachelor_thesis_headline_clist}
% 本科生模板的页眉标题与封面大标题常量列表，其中封面大标题不适用于硕士、博士学位论文，只适用于本科生毕业设计（论文）及其衍生物。
%    \begin{macrocode}
\clist_const:Nn \c_@@_bachelor_thesis_header_clist
  {
    上海电力大学本科生毕业设计（论文）,
    上海电力大学硕士学位论文,
    上海电力大学博士学位论文,
  }
\clist_const:Nn \c_@@_bachelor_thesis_headline_clist
  {
    本科生毕业设计（论文）,
  }
%    \end{macrocode}
% \end{variable}
%
% \subsubsection{l3keys 接口键值对定义}
%
% 定义 |suepthesis| 键值对类。
%    \begin{macrocode}
\keys_define:nn { suepthesis }
{
  info .meta:nn = { suepthesis / info } {#1},
  misc .meta:nn = { suepthesis / misc } {#1},
  cover .meta:nn = { suepthesis / cover } {#1},
  style .meta:nn = { suepthesis / style } {#1},
  option .meta:nn = { suepthesis / option } {#1},
  TOC .meta:nn = { suepthesis / TOC } {#1},
  appendices .meta:nn = { suepthesis / appendices } {#1},
  publications .meta:nn = { suepthesis / publications } {#1},
  const .meta:nn = { suepthesis / const } {#1},
}
%    \end{macrocode}
%
% 定义 |suepthesis/option| 键值对类。
%    \begin{macrocode}
\keys_define:nn { suepthesis / option }
{
  type .choice:,
  type .value_required:n = true,
  type .choices:Vn =
    \c_@@_thesis_type_clist
    {
      \int_set_eq:NN \g_@@_thesis_type_int \l_keys_choice_int
    },
  type .initial:n = master,
  degree .choice:,
  degree .value_required:n = true,
  degree .choices:Vn =
    \c_@@_thesis_degree_type_clist
    {
      \int_set_eq:NN \g_@@_thesis_degree_type_int \l_keys_choice_int
    },
  degree .initial:n = academic,
  % 在双面打印模式下，需要放在奇数页的页面后会自动插入一个空白页，以方便直接在打印机上“双面打印”。
  twoside .bool_gset:N = \g_@@_twoside_bool,
  twoside .initial:n = true,
  blindPeerReview .bool_gset:N = \g_@@_blind_mode_bool,
  twoside .initial:n = false,
  ctex .tl_set:N = \l_@@_options_to_ctex_tl,
  quirks .bool_gset:N = \g_@@_quirks_mode_bool,
  % xeCJK
  autoFakeBold .tl_set:N = \g_@@_auto_fake_bold_tl,
  autoFakeBold .initial:n = {3},
}
%    \end{macrocode}
%
% 定义 |suepthesis/cover| 键值对类。
%    \begin{macrocode}
\keys_define:nn { suepthesis / cover }
  {
    date .tl_set:N = \l_@@_cover_date_tl,
    headerImage .tl_set:N = \l_bit_coverheaderimage_tl,
    xiheiFont .tl_set:N = \l_@@_cover_xihei_font_path_tl,
    xiheiFont .default:n = {STXihei},
    %% cover entry
    delimiter .tl_set:N = \l_@@_cover_delimiter_tl,
    labelAlign .tl_set:N = \l_@@_cover_label_align_tl,
    labelAlign .initial:n = {r},
    valueAlign .tl_set:N = \l_@@_cover_value_align_tl,
    valueAlign .initial:n = {c},
    labelMaxWidth .dim_set:N = \l_@@_cover_label_max_width_dim,
    valueMaxWidth .dim_set:N = \l_@@_cover_value_max_width_dim,
    autoWidthPadding .dim_set:N = \l_@@_cover_auto_width_padding_dim,
    autoWidthPadding .initial:n = {0.25em},
    autoWidth .bool_set:N = \l_@@_cover_auto_width_bool,
    autoWidth .initial:n = {true},
    underlineThickness .dim_set:N = \l_@@_cover_underline_thickness_dim,
    underlineThickness .initial:n = {1pt},
    underlineOffset .dim_set:N = \l_@@_cover_underline_offset_dim,
    underlineOffset .initial:n = { -10pt },
    hideCoverInPeerReview .bool_set:N = \l_@@_cover_hide_cover_in_peer_review_bool,
    % 此处暂且填充默认值为`false`，待确定`\g_@@_thesis_type_int`后再根据论文类型覆盖默认值
    hideCoverInPeerReview .initial:n = {false},
    % 研究生的「特殊类型」
    showSpecialTypeBox .bool_set:N = \l_@@_cover_show_special_type_box_bool,
    showSpecialTypeBox .initial:n = {false},
    % 本科英文模板也可以使用中文封面
    prefer-zh .bool_set:N = \l_@@_cover_prefer_zh_bool,
    prefer-zh .initial:n = {false},
    % 本科英文模板使用中文封面时可能需要调换中英文标题顺序
    reverse-titles .bool_set:N = \l_@@_cover_reverse_titles_bool,
    reverse-titles .initial:n = {false},
    % 本科英文模板可加入中文标题
    addTitleZh .bool_set:N = \l_@@_cover_add_titlezh_bool,
    addTitleZh .initial:n = {true},
  }
%    \end{macrocode}
%
% 定义 |suepthesis/info| 键值对类。
%    \begin{macrocode}
\keys_define:nn { suepthesis / info }
  {
    % 中文题目
    title .tl_set:N = \l_@@_value_title_tl,
    title .initial:n = {形状记忆聚氨酯的合成及其在织物中的应用（示例）},
    % 英文题目
    titleEn .tl_set:N = \l_@@_value_title_en_tl,
    titleEn .initial:n = {Synthesis~and~Application~on~Texttiles~of~the~Shape~Memory~Polyurethane~(example)},

    institute .tl_set:N = \l_@@_value_institute_tl,
    major .tl_set:N = \l_@@_value_major_tl,
    author .tl_set:N = \l_@@_value_author_tl,
    studentId .tl_set:N = \l_@@_value_student_id_tl,
    % 指导教师
    supervisor .tl_set:N = \l_@@_value_supervisor_tl,

    keywords .tl_set:N = \l_@@_value_keywords_tl,
    keywords .initial:n = {聚氨酯,形状记忆,织物（示例）},
    keywordsEn .tl_set:N = \l_@@_value_keywords_en_tl,
    keywordsEn .initial:n = {Polyurethane,Shape Memory,Textiles (example)},

    % 中图分类号，研究生学位论文使用
    classification .tl_set:N = \l_@@_value_classification_tl,
    classification .initial:n = {TQ~028.1（示例）},
    degree .tl_set:N = \l_@@_value_degree_class_tl,
    finalizationDate .tl_set:N = \l_@@_value_finalization_date_tl,
  }
%    \end{macrocode}
%
% 定义 |suepthesis/misc| 键值对类。
%    \begin{macrocode}
\keys_define:nn { suepthesis / misc }
  {
    % 表格字体大小，默认为 5 号字体。
    tabularFontSize .tl_set:N = \l_@@_misc_tabular_font_size_tl,
    tabularFontSize .initial:n = {5},
    arialFont .tl_set:N = \l_@@_misc_arial_font_path_tl,
    autoref / algo .code:n = {
      \AtBeginDocument{
        % 定义算法标题
        % 针对 algorithm 宏包
        \tl_set:Nn \ALG@name {#1}
        % 针对 algorithm2e 宏包
        \tl_set:Nn \algorithmcfname {#1}

        % 定义算法的 autoref
        % algorithm2e 宏包会覆写它，所以我们必须AtBeginDocument时再修改
        \tl_set:Nn \algorithmautorefname {#1}
      }
    },
    autoref / algo .initial:n = {\g_@@_const_autoref_algo_tl},
    autoref / them .tl_set:N = \themautorefname,
    autoref / them .initial:n = {\g_@@_const_autoref_them_tl},
    autoref / lem .tl_set:N = \lemautorefname,
    autoref / lem .initial:n = {\g_@@_const_autoref_lem_tl},
    autoref / prop .tl_set:N = \propautorefname,
    autoref / prop .initial:n = {\g_@@_const_autoref_prop_tl},
    autoref / cor .tl_set:N = \corautorefname,
    autoref / cor .initial:n = {\g_@@_const_autoref_cor_tl},
    autoref / axi .tl_set:N = \axiautorefname,
    autoref / axi .initial:n = {\g_@@_const_autoref_axi_tl},
    autoref / defn .tl_set:N = \defnautorefname,
    autoref / defn .initial:n = {\g_@@_const_autoref_defn_tl},
    autoref / conj .tl_set:N = \conjautorefname,
    autoref / conj .initial:n = {\g_@@_const_autoref_conj_tl},
    autoref / exmp .tl_set:N = \exmpautorefname,
    autoref / exmp .initial:n = {\g_@@_const_autoref_exmp_tl},
    autoref / case .tl_set:N = \caseautorefname,
    autoref / case .initial:n = {\g_@@_const_autoref_case_tl},
    autoref / rem .tl_set:N = \remautorefname,
    autoref / rem .initial:n = {\g_@@_const_autoref_rem_tl},
    hideLinks .bool_set:N = \l_@@_misc_hide_links_bool,
    hideLinks .initial:n = {true},
    autoref / figure .tl_set:N = \figureautorefname,
    autoref / figure .initial:n = {\g_@@_const_autoref_fig_tl},
    autoref / table .tl_set:N = \tableautorefname,
    autoref / table .initial:n = {\g_@@_const_autoref_tab_tl},
    autoref / equ .tl_set:N = \equationautorefname,
    autoref / equ .initial:n = {\g_@@_const_autoref_equ_tl},
    % 浮动体相关的各种间距
    floatSeparation .tl_set:N = \l_@@_misc_float_separation_tl,
    floatSeparation .initial:n = {0},
    algorithmSeparation .tl_set:N = \l_@@_misc_algorithm_separation_tl,
    algorithmSeparation .initial:n = {12pt plus 4pt minus 4pt},
    tabularRowSeparation .tl_set:N = \l_@@_misc_tabular_row_separation_tl,
    tabularRowSeparation .initial:n = {1},
  }
%    \end{macrocode}
%
% 定义 |suepthesis/const| 键值对类。
%    \begin{macrocode}
\keys_define:nn { suepthesis / const }
  {
    autoref .meta:nn = { suepthesis / const / autoref } { #1 },
    style .meta:nn = { suepthesis / const / style } { #1 },
    info .meta:nn = { suepthesis / const / info } { #1 },
    heading .meta:nn = { suepthesis / const / heading } { #1 },
  }
\keys_define:nn { suepthesis / const / autoref }
  {
    algo .tl_set:N = \g_@@_const_autoref_algo_tl,
    algo .initial:n = {\@@_get_const:N {algo}},
    them .tl_set:N = \g_@@_const_autoref_them_tl,
    them .initial:n = {\@@_get_const:N {them}},
    lem .tl_set:N = \g_@@_const_autoref_lem_tl,
    lem .initial:n = {\@@_get_const:N {lem}},
    prop .tl_set:N = \g_@@_const_autoref_prop_tl,
    prop .initial:n = {\@@_get_const:N {prop}},
    cor .tl_set:N = \g_@@_const_autoref_cor_tl,
    cor .initial:n = {\@@_get_const:N {cor}},
    axi .tl_set:N = \g_@@_const_autoref_axi_tl,
    axi .initial:n = {\@@_get_const:N {axi}},
    defn .tl_set:N = \g_@@_const_autoref_defn_tl,
    defn .initial:n = {\@@_get_const:N {defn}},
    conj .tl_set:N = \g_@@_const_autoref_conj_tl,
    conj .initial:n = {\@@_get_const:N {conj}},
    exmp .tl_set:N = \g_@@_const_autoref_exmp_tl,
    exmp .initial:n = {\@@_get_const:N {exmp}},
    case .tl_set:N = \g_@@_const_autoref_case_tl,
    case .initial:n = {\@@_get_const:N {case}},
    rem .tl_set:N = \g_@@_const_autoref_rem_tl,
    rem .initial:n = {\@@_get_const:N {rem}},
    figure .tl_set:N = \g_@@_const_autoref_fig_tl,
    figure .initial:n =  {\@@_get_const:N {fig}},
    table .tl_set:N = \g_@@_const_autoref_tab_tl,
    table .initial:n =  {\@@_get_const:N {tab}},
    equ .tl_set:N = \g_@@_const_autoref_equ_tl,
    equ .initial:n = {\@@_get_const:N {equ}},
  }
\keys_define:nn { suepthesis / const / style }
  {
    substituteSymbol .tl_set:N = \g_@@_const_substitute_symbol_tl,
    substituteSymbol .initial:n = {*},
  }
\keys_define:nn { suepthesis / const / info }
  {
    degree .tl_set:N = \g_@@_const_info_degree_tl,
    degree .initial:n = {
      \c_@@_graduate_label_degree_tl
    },
    major .tl_set:N = \g_@@_const_info_major_tl,
    major .initial:n = {
      \@@_if_graduate:TF {
        \c_@@_graduate_label_major_tl
      } {
        \@@_get_const:N {major}
      }
    },
  }
\keys_define:nn { suepthesis / const / heading }
  {
    acknowledgements .tl_set:N = \g_@@_const_heading_acknowledgements_tl,
    acknowledgements .initial:n = {
      \@@_get_const:N {ack}
    },
  }
%    \end{macrocode}
%
% 定义 |suepthesis/style| 键值对类。
%    \begin{macrocode}
\keys_define:nn { suepthesis / style }
{
  head .tl_set:N = \l_@@_style_head_tl,
  head .initial:n = {
    \clist_item:Nn \c_@@_bachelor_thesis_header_clist \g_@@_thesis_type_int
  },
  headline .tl_set:N = \l_@@_style_headline_tl,
  headline .initial:n = {
    \clist_item:Nn \c_@@_bachelor_thesis_headline_clist \g_@@_thesis_type_int
  },
  bibliographyIndent .bool_set:N = \l_@@_style_bibliography_indent_bool,
  bibliographyIndent .initial:n = {true},
  pageVerticalAlign .choices:nn = {top, scattered} {
    \tl_if_eq:NnTF \l_keys_choice_tl {top}
      { \raggedbottom }
      { \flushbottom }
  },
  pageVerticalAlign .initial:n = {top},
  non-CJK-font-in-headings .choice:,
  non-CJK-font-in-headings / serif .code:n = { \bool_set_false:N \l_@@_arial_as_title_font_bool },
  non-CJK-font-in-headings / sans .code:n = { \bool_set_true:N \l_@@_arial_as_title_font_bool },
  non-CJK-font-in-headings .initial:n = {serif},
  % 数学字体配置
  mathFont .choices:nn = {
    asana, bonum, cm, concrete, dejavu, erewhon, euler,
    fira, garamond, gfsneohellenic, kp, libertinus, lm, newcm,
    pagella, schola, stix, stix2, termes, xcharter, xits, none,
    } { \tl_set_eq:NN \l_@@_style_math_font_tl \l_keys_choice_tl },
  mathFont .initial:n = {cm},
  % Options that will be pass to `unicode-math` pkgs.
  unicodeMathOptions .tl_set:N = \l_@@_unicode_math_options_tl,
  % Windows 平台开启宋体伪粗体。
  windowsSimSunFakeBold .bool_set:N = \l_@@_style_windows_simsum_fake_bold,
  windowsSimSunFakeBold .initial:n = {false},
  % 控制英文是否使用 hyphen 进行换行
  hyphen .bool_set:N = \l_@@_style_hyphen_bool,
  hyphen .initial:n = {true},
  % 控制公式和上下文的距离
  mathAboveDisplaySkip .dim_set:N = \l_@@_style_math_above_display_skip_dim,
  mathAboveDisplaySkip .initial:n = {10pt},
  mathBelowDisplaySkip .dim_set:N = \l_@@_style_math_below_display_skip_dim,
  mathBelowDisplaySkip .initial:n = {10pt},
  betterTimesNewRoman .bool_set:N = \l_@@_style_better_new_roman_bool,
  betterTimesNewRoman .initial:n = {false},
}
%    \end{macrocode}
%
% 定义 |suepthesis/TOC| 键值对类。
%    \begin{macrocode}
\keys_define:nn { suepthesis / TOC }
{
  title .tl_set:N = \l_@@_toc_title_tl,
  title .initial:n = {
    \@@_get_const:N {toc}
  },
  abstract .bool_set:N = \l_@@_add_abstract_to_toc_bool,
  abstract .initial:n = {true},
  abstractEn .bool_set:N = \l_@@_add_abstract_en_to_toc_bool,
  abstractEn .initial:n = {true},
  TOC .bool_set:N = \l_@@_add_toc_to_toc_bool,
  TOC .initial:n = {false},
  symbols .bool_set:N = \l_@@_add_symbols_to_toc_bool,
  symbols .initial:n = {true},
}
%    \end{macrocode}
%
% 定义 |suepthesis/appendices| 键值对类。
%    \begin{macrocode}
\keys_define:nn { suepthesis / appendices }
{
  chapterLevel .bool_set:N = \l_@@_appendices_chapter_level_bool,
  title .tl_set:N = \l_@@_appendices_title_tl,
  TOCTitle .tl_set:N = \l_@@_appendix_toc_title_tl,
}
%    \end{macrocode}
%
% 定义 |suepthesis/appendices| 键值对类。
%    \begin{macrocode}
\keys_define:nn { suepthesis / publications }
{
  % mode .choice:,
  % mode .value_required:n = true,
  % mode .choices:Vn =
  %   \c_@@_publication_modes_clist
  %   {
  %     \int_new:N \l_@@_publication_mode_int
  %     \int_set:Nn \l_@@_publication_mode_int \l_keys_choice_int
  %   },
  % mode .initial:n = biblatex,
  sorting .bool_set:N = \l_@@_publications_sorting_bool,
  sorting .initial:n = {true},
  omit .bool_set:N = \l_@@_publications_omit_bool,
  omit .initial:n = {false},
  maxbibnames .int_set:N = \l_@@_publications_maxbibnames_int,
  maxbibnames .initial:n = {10},
  minbibnames .int_set:N = \l_@@_publications_minbibnames_int,
  minbibnames .initial:n = {10},
}
%    \end{macrocode}
% 在宏加载时，处理 |suepthesis/option| 中的值。使得 |suepthesis|
% 宏包的模板选项可以在宏加载时生效。
%    \begin{macrocode}
\ProcessKeysOptions { suepthesis / option }
%    \end{macrocode}
% 确定 |suepthesis/option| 中的 |\g_@@_thesis_type_int| 后，根据论文类型自动覆盖某些选项的默认值。
%    \begin{macrocode}
\@@_if_graduate:TF {
  \keys_set:nn {suepthesis} {cover/hideCoverInPeerReview = false}
} {
  \keys_set:nn {suepthesis} {cover/hideCoverInPeerReview = true}
}
%    \end{macrocode}
%
% \subsubsection{处理模板选项}
%
% 如果设置 blindPeerReview 选项，则抑制 twoside 选项。
%    \begin{macrocode}
\bool_if:NT \g_@@_blind_mode_bool {
  \bool_set_false:N \g_@@_twoside_bool
}
%    \end{macrocode}
%
% 如果没有开启双面打印选项，则在 ctexbook 中开启单面打印选项。
% 允许 chapter 直接另起一页（即使是偶数（左手）页）。
%    \begin{macrocode}
\bool_if:NTF \g_@@_twoside_bool {} {
  \PassOptionsToClass{oneside}{ctexbook}
}
\PassOptionsToClass{openany}{ctexbook}
%    \end{macrocode}
%
% 将 |suepthesis/option/ctex| 中的值传递给 ctexbook 模板类。
%    \begin{macrocode}
% Any extra option passed by user will be passed to ctexbook.
\DeclareOption*{
  \PassOptionsToClass{\l_@@_options_to_ctex_tl}{ctexbook}
}
%    \end{macrocode}
%
% 抑制 fontspec 宏包关于字体的警告信息。
% 手动开启伪粗体、伪斜体。
%    \begin{macrocode}
\PassOptionsToPackage{quiet,AutoFakeBold=\g_@@_auto_fake_bold_tl,AutoFakeSlant}{xeCJK}
%    \end{macrocode}
% 加载 ctexbook 模板类。
%    \begin{macrocode}
\ProcessOptions\relax
\LoadClass[zihao=-4,]{ctexbook}
%    \end{macrocode}
%
% \subsubsection{定义模板类样式}
%
% 加载所需的宏包。
%    \begin{macrocode}
\RequirePackage{geometry}
\RequirePackage[table,xcdraw]{xcolor}
\RequirePackage{xeCJK}
\RequirePackage {xeCJKfntef}
% 恢复数学行距（restoremathleading），同时避免改变正文行距。
% （ctex 默认 linespread 1.3 × LaTeX 默认倍数 1.2 = 1.56）
\RequirePackage[bodytextleadingratio=1.56]{zhlineskip}
\RequirePackage{titletoc}
\RequirePackage{graphicx}
\RequirePackage{fancyhdr}
\RequirePackage{pdfpages}
\RequirePackage[nodisplayskipstretch]{setspace}
\RequirePackage{booktabs}
\RequirePackage{multirow}
\RequirePackage{tikz}
\RequirePackage{etoolbox}
% Hide color and border in hyperref.
\RequirePackage[bookmarksnumbered]{hyperref}
% 详见 `caption` 宏包手册和
% https://github.com/CTeX-org/forum/issues/86
\RequirePackage[strut=off]{caption}
\RequirePackage{array}
\RequirePackage{amsmath}
\RequirePackage{amssymb}
\RequirePackage{pifont}
\RequirePackage{amsthm}
\RequirePackage{pdfpages}
\RequirePackage{listings}
\RequirePackage{enumitem}
\RequirePackage{fmtcount}
%    \end{macrocode}
%
% 抑制 \pkg{hyperref} 中对 |\hskip| 的 warning 信息。
%    \begin{macrocode}
\pdfstringdefDisableCommands{%
  \let\quad\empty
}
%    \end{macrocode}
%
% 设置页眉字号，页边距。
%
% 需要注意的是，根据 \pkg{geometry} 的规则，
% |headsep| 和 |footskip| 分别受到 |top| 与 |bottom| 的影响。
% 所以你能看到在计算 |headsep| 与 |footskip| 时，我们
% 首先计算了相应的偏移量。
%    \begin{macrocode}
\@@_if_graduate:TF {
  \geometry{
    a4paper,
    left=2.6cm,
    bottom=2.5cm + 7bp,
    top=3cm + 7bp,
    right=2.6cm,
    % `headsep' is affected by `top' option.
    headsep = 3cm + 7bp - 2.5cm - 15bp,
    headheight = 15 bp,
    % `footskip' is affected by `bottom' option.
    footskip = 2.5cm + 7bp - 1.8cm,
  }
} {
  \geometry{
    a4paper,
    left=3cm,
    bottom=2.6cm + 7bp,
    top=3.5cm + 7bp,
    right=2.6cm,
    % `headsep' is affected by `top' option.
    headsep = 3.5cm + 7bp - 2.4cm - 20bp,
    headheight = 20 bp,
    % `footskip' is affected by `bottom' option.
    footskip = 2.6cm + 7bp - 2cm,
  }
}
%    \end{macrocode}
%
%  \subsubsection{定义字体相关选项}
%
% 设置 Times New Roman 字体。
% 根据学校规范要求，默认情况下也使用 Times New Roman 字体。
%    \begin{macrocode}
\ctex_at_end_preamble:n {
  \bool_if:NTF \l_@@_style_better_new_roman_bool {
    \defaultfontfeatures[TeX~Gyre~Termes]
    {
     Extension      = .otf ,
     UprightFont    = texgyretermes-regular,
     BoldFont       = texgyretermes-bold,
     ItalicFont     = texgyretermes-italic,
     BoldItalicFont = texgyretermes-bolditalic,
    }
    \setmainfont{TeX~Gyre~Termes}
  }{
    \setmainfont{Times~New~Roman}
    \setromanfont{Times~New~Roman}
  }
}
%    \end{macrocode}
%
% \begin{macro}{\@@_font_path:}
% 当选择使用字体文件配置字体时，设置字体文件路径。
%    \begin{macrocode}
\cs_new:Npn \@@_font_path:
  {
    \str_if_eq:NNTF { \l_@@_font_type_tl } { font }
      { }
      { Path = \l_@@_font_path_tl / , }
  }
%    \end{macrocode}
% \end{macro}
%
% \begin{macro}{\@@_load_unicode_math_pkg:}
% 加载\pkg{unicode-math}宏包。
%    \begin{macrocode}
\cs_new:Npn \@@_load_unicode_math_pkg:
  {
    \PassOptionsToPackage { \l_@@_unicode_math_options_tl } { unicode-math }
    \RequirePackage { unicode-math }
  }
%    \end{macrocode}
% \end{macro}
%
% \begin{macro}{\@@_define_math_font:nn}
% 批量定义数学字体配置。
% \begin{arguments}
%   \item 配置名称。
%   \item 字体名称。
% \end{arguments}
%    \begin{macrocode}
\cs_new:Npn \@@_define_math_font:nn #1#2
  {
    \cs_new:cpn { @@_load_math_font_ #1 : }
      {
        \@@_load_unicode_math_pkg:
        \setmathfont { #2 }
      }
  }
\clist_map_inline:nn
  {
    { asana          } { Asana-Math.otf             },
    { concrete       } { Concrete-Math.otf          },
    { erewhon        } { Erewhon-Math.otf           },
    { euler          } { Euler-Math.otf             },
    { fira           } { FiraMath-Regular.otf       },
    { garamond       } { Garamond-Math.otf          },
    { gfsneohellenic } { GFSNeohellenicMath.otf     },
    { kp             } { KpMath-Regular.otf         },
    { libertinus     } { LibertinusMath-Regular.otf },
    { lm             } { latinmodern-math.otf       },
    { newcm          } { NewCMMath-Regular.otf      },
    { stix           } { STIXMath-Regular.otf       },
    { stix2          } { STIXTwoMath-Regular.otf    },
    { xcharter       } { XCharter-Math.otf          },
    { xits           } { XITSMath-Regular.otf       },
    { bonum          } { texgyrebonum-math.otf      },
    { dejavu         } { texgyredejavu-math.otf     },
    { pagella        } { texgyrepagella-math.otf    },
    { schola         } { texgyreschola-math.otf     },
    { termes         } { texgyretermes-math.otf     }
  }
  { \@@_define_math_font:nn #1 }
%    \end{macrocode}
% \end{macro}
%
% \begin{macro}{\@@_load_math_font_cm:}
% 数学字体配置 |cm|。
%    \begin{macrocode}
\cs_new:Npn \@@_load_math_font_cm: { }
%    \end{macrocode}
% \end{macro}
%
% \begin{macro}{\@@_load_math_font_none:}
% 数学字体配置 |none|。
%    \begin{macrocode}
\cs_new:Npn \@@_load_math_font_none: { }
%    \end{macrocode}
% \end{macro}
%
% \begin{macro}{\@@_load_font:}
% 加载数学字体
%    \begin{macrocode}
\cs_new:Npn \@@_load_font:
  {
    \use:c { @@_load_math_font_  \l_@@_style_math_font_tl  : }
  }
%    \end{macrocode}
% \end{macro}
%
%  \paragraph{定义导言区末尾加载内容}
%
% 在 |preamble| 中，加载各个模板需要的字体。
%    \begin{macrocode}
\ctex_at_end_preamble:n {
  % 针对 Windows 字体采用 Fake Bold 宋体
  \bool_if:NT \l_@@_style_windows_simsum_fake_bold
  {
    \RequirePackage{ifplatform}
    \ifwindows
      \setCJKmainfont{SimSun}[AutoFakeBold,AutoFakeSlant]
    \fi
  }

  % 在导言区末尾加载数学字体。
  \@@_load_font:

  % misc / hideLinks 选项
  \bool_if:NTF \l_@@_misc_hide_links_bool
    {
      \hypersetup { hidelinks }
    } {
      \definecolor{blue}{RGB}{10,10,110}
      \hypersetup{
        colorlinks=true,
      }
    }

  \bool_if:NTF \l_@@_arial_as_title_font_bool {
    % 手动指定时要加载 Arial
    \newfontfamily\arialfamily{Arial}
  } {
    % 即使未指定，本科全英文专业模板的声明也需要 Arial
    \@@_if_thesis_int_type:nT {3} {
      \newfontfamily\arialfamily{Arial}
    }
  }
  % 无论中英文，封面都可能需要细黑体。
  \tl_if_blank:VTF \l_@@_cover_xihei_font_path_tl {}
  {
    \setCJKfamilyfont{xihei}[AutoFakeBold,AutoFakeSlant]
      {\l_@@_cover_xihei_font_path_tl}
  }

  % 对于本科全英文专业模板，需要自定义日期格式。
  \@@_if_thesis_int_type:nT {3} {
    \RequirePackage[en-US]{datetime2}
    \RequirePackage{indentfirst}
    \DTMlangsetup[en-US]{dayyearsep={\space}}
  }

  % Define biblatex category if it was imported.
  % 这部分是给研究生模板中的
  % 「攻读学位期间发表论文与研究成果清单」使用的。
  \cs_if_exist:NT \DeclareBibliographyCategory {
    \DeclareBibliographyCategory{mypub}
  }

  % Define biblatex strings if it was imported.
  % 这部分是给研究生模板中的
  % gbpunctin = false 时使用的。
  \cs_if_exist:NT \DefineBibliographyStrings {
    \DefineBibliographyStrings{english}{in={}}
    \DefineBibliographyStrings{english}{incn={}}
  }

  % 修改 biblatex 中「专利」（patent）部分的著录格式。
  % 主要根据北理工自定义的规范，参考 biblatex 和
  % biblatex-gb7714-2015 的实现修改而成。
  %
  % 默认不开启，因为此修改可能会产生其他边界问题。
  \bool_if:NT \g_@@_quirks_mode_bool {
    \cs_if_exist:NT \DeclareBibliographyDriver {
      %
      %   重设专利title的输出，将文献类型标识符输出出去
      %
      \newbibmacro*{patenttitle}{%原输出来自biblatex.def文件
        \ifboolexpr{%
          test{\iffieldundef{title}}%
          and%
          test{\iffieldundef{subtitle}}%
        }%
          {}%
          {\printtext[title]{\bibtitlefont%
             \printfield[titlecase]{title}%
             \ifboolexpr{test {\iffieldundef{subtitle}}}%这里增加了对子标题的判断，解决不判断多一个点的问题
             {}{\setunit{\subtitlepunct}%
             \printfield[titlecase]{subtitle}}%
              \iftoggle{bbx:gbtype}{\printfield[gbtypeflag]{usera}}{}%
             \iffieldundef{titleaddon}{}%判断一下titleaddon，否则直接加可能多一个空格
              {\setunit{\subtitlepunct}\printfield{titleaddon}}%
              % ：地区
              \setunit{\subtitlepunct}\iflistundef{location}
                {}
                {\setunit*{\subtitlepunct}%
                 \printtext{%[parens]
                   \printlist[][-\value{listtotal}]{location}}}%
              % ，专利号
              \setunit{\addcomma\addspace}\printfield{number}%写专利号
              \setunit{\addcomma\addspace}
              \usebibmacro{newsdate}%
           }%
        }%
      }

      %
      %   重定义专利文献驱动
      %
        \DeclareBibliographyDriver{patent}{%源来自standard.BBX
        \usebibmacro{bibindex}%
        \usebibmacro{begentry}%
        \usebibmacro{author}%
        \ifnameundef{author}{}{\setunit{\labelnamepunct}\newblock}%这一段用于去除作者不存在时多出的标点
        \usebibmacro{patenttitle}%给出专利专用的标题输出
        \iftoggle{bbx:gbstrict}{}{%
          \newunit%
          \printlist{language}%
          \newunit\newblock
          \usebibmacro{byauthor}
        }%
        \newunit\newblock
        \printfield{type}%
        \setunit*{\addspace}%
        \newunit\newblock
        \usebibmacro{byholder}%
        \newunit\newblock
        \printfield{note}%
        \newunit\newblock
        \usebibmacro{doi+eprint+url}%
        \newunit\newblock
        \usebibmacro{addendum+pubstate}%
        \setunit{\bibpagerefpunct}\newblock
        \usebibmacro{pageref}%
        \newunit\newblock
        \iftoggle{bbx:related}
          {\usebibmacro{related:init}%
           \usebibmacro{related}}
          {}%
        \usebibmacro{annotation}\usebibmacro{finentry}}
    }
  }
}
%    \end{macrocode}
%
% \begin{macro}{\xihei:n}
% 定义细黑字体。
%    \begin{macrocode}
\cs_new:Npn \xihei:n #1 {
  \xeCJK_family_if_exist:nTF {xihei} {
    \CJKfamily{xihei} #1
  }{
    \heiti #1
  }
}
%    \end{macrocode}
% \end{macro}
%
% \begin{macro}{\l_@@_unnumchapter_style_cs:n}
% 定义无序章节的样式。
%    \begin{macrocode}
\cs_new:Npn \l_@@_unnumchapter_style_cs:n #1 {
  % 本科全英文、研究生学位论文需要加粗
  \int_compare:nNnTF {\g_@@_thesis_type_int} > {2}
  {
    \bfseries #1
  } {
    \mdseries #1
  }
}
%    \end{macrocode}
% \end{macro}
%
% 定义 \pkg{fancyhdr} 的页眉页脚。
%    \begin{macrocode}
\fancypagestyle{suepthesis}{
  \fancyhf{}
  % 定义页眉、页码
  \fancyhead[C]{
    \zihao{\int_use:N \g_@@_head_zihao_int}
    \ziju{0.08}
    \songti{\tl_use:N \l_@@_style_head_tl}
  }
  \fancyfoot[C]{\songti\zihao{5} \thepage}
  % 页眉分割线稍微粗一些
  \RenewDocumentCommand \headrulewidth {} {0.6pt}
}
%    \end{macrocode}
%
%  定义 \pkg{ctex} 的章节标题形式。
%    \begin{macrocode}
\ctexset{chapter={
    number = {\arabic{chapter}},
    format = {\heiti \bfseries \centering \zihao{3}},
    nameformat = {},
    titleformat = {},
    aftername = \hspace{9bp},
    pagestyle = suepthesis,
    beforeskip = 8bp,
    afterskip = 32bp,
    fixskip = true,
    lofskip = 0cm,
    lotskip = 0cm,
  }
}

\ctexset{section={
    number = {\thechapter.\hspace{1bp}\arabic{section}},
    format = {\heiti \raggedright \bfseries \zihao{-3}},
    nameformat = {},
    titleformat = {},
    aftername = \hspace{8bp},
    beforeskip = 20bp,
    afterskip = 17bp,
    fixskip = true,
  }
}

\ctexset{subsection={
    number = {
        \thechapter.\hspace{1bp}
        \arabic{section}.\hspace{1bp}
        \arabic{subsection}
    },
    format = { \heiti \bfseries \raggedright \zihao{-4}},
    nameformat = {},
    titleformat = {},
    aftername = \hspace{7bp},
    beforeskip = 17bp,
    afterskip = 17bp,
    fixskip = true,
  }
}

\ctexset{
  secnumdepth = 3,
  subsubsection={
    numbering = true,
    number = {
        \arabic{chapter}.\hspace{1bp}
        \arabic{section}.\hspace{1bp}
        \arabic{subsection}.\hspace{1bp}
        \arabic{subsubsection}
    },
    format={\heiti \raggedright \zihao{-4}},
    nameformat = {},
    titleformat = {},
    beforeskip=14bp,
    afterskip=14bp,
    fixskip=true,
  }
}
%    \end{macrocode}
%
%  定义 TOC 样式。
%    \begin{macrocode}
\addtocontents{toc}{\protect\hypersetup{hidelinks}}

\@@_if_graduate:TF {
  % 对于研究生模板，定义各章标题为宋体四号。
  \titlecontents{chapter}[0pt]{\songti \zihao{-4}}
  {\thecontentslabel\hspace{\ccwd}}{}
  {\hspace{.5em}\titlerule*{.}\contentspage}
  % section 标题为宋体小四号。缩进为两个字符宽度。
  \titlecontents{section}[2\ccwd]{\songti \zihao{-4}}
  {\thecontentslabel\hspace{\ccwd}}{}
  {\hspace{.5em}\titlerule*{.}\contentspage}
  % subsection 标题为宋体小四号。缩进为四个字符宽度。
  \titlecontents{subsection}[4\ccwd]{\songti \zihao{-4}}
  {\thecontentslabel\hspace{\ccwd}}{}
  {\hspace{.5em}\titlerule*{.}\contentspage}
} {
  % 对于其他，定义各章标题为宋体小四号。
  \titlecontents{chapter}[0pt]{\songti \zihao{-4}}
  {\thecontentslabel\hspace{\ccwd}}{}
  {\hspace{.5em}\titlerule*{.}\contentspage}
  % section 标题为宋体小四号。
  \titlecontents{section}[1\ccwd]{\songti \zihao{-4}}
  {\thecontentslabel\hspace{\ccwd}}{}
  {\hspace{.5em}\titlerule*{.}\contentspage}
  % subsection 标题为宋体小四号。
  \titlecontents{subsection}[2\ccwd]{\songti \zihao{-4}}
  {\thecontentslabel\hspace{\ccwd}}{}
  {\hspace{.5em}\titlerule*{.}\contentspage}
}
% listoffigure 样式优化
\titlecontents{figure}[0pt]{\songti\zihao{-4}}
    {\figurename~\thecontentslabel\quad}{\hspace*{-1.5cm}}
    {\hspace{.5em}\titlerule*{.}\contentspage}
% listoftable 样式优化
\titlecontents{table}[0pt]{\songti\zihao{-4}}
    {\tablename~\thecontentslabel\quad}{\hspace*{-1.5cm}}
    {\hspace{.5em}\titlerule*{.}\contentspage}
%    \end{macrocode}
%
%  \subsubsection{定义样式相关函数}
%
% \begin{macro}{\frontmatter}
% 定义前置内容的页面样式。
%    \begin{macrocode}
\RenewDocumentCommand \frontmatter {} {
  \pagenumbering{Roman}
  % 这部分的章节标题不进行编号。
  \ctexset{
    chapter = {
      numbering = false,
    }
  }
  \linespread{1.53}\selectfont
  \pagestyle{suepthesis}

  % 调整表格内容字号（默认五号）和各行之间的距离。
  %
  % 由于这种方式会影响所有的表格，
  % 所以我们尽可能延迟这种影响。
  %
  % 不过，在目前的代码实现中没有在封面
  % 之类的地方使用表格，所以目前即使放在
  % preamble 中也不会有影响。
  %
  % 支持标准tabular、tabular*环境和宏包tabularx、longtable。
  %
  % 为保证各种表格效果一致，要先手动重置setspace宏包漏掉的longtable；
  \AtBeginEnvironment {longtable} {\singlespacing}
  % 之后再统一设置。
  \clist_map_inline:nn
    {tabular, tabular*, tabularx, longtable}
    {
      \AtBeginEnvironment {##1} {
        % 字号只想设置表格内容，不想影响caption。
        % 一般caption在环境之外，自然不受影响；
        % 而longtable的caption虽在环境内，但caption宏包能正常处理。
        \zihao{\l_@@_misc_tabular_font_size_tl}
        % 各行间距只想影响表格，不想影响矩阵，因此也必须在钩子中设置。
        \cs_set:Npn \arraystretch {\l_@@_misc_tabular_row_separation_tl}
      }
    }
}
%    \end{macrocode}
% \end{macro}
%
% \begin{macro}{\mainmatter}
% 主体内容的页面样式。
%    \begin{macrocode}
\RenewDocumentCommand \mainmatter {} {
  % 另起一个空页，以便于后续的章节标题编号。
  % \clearpage
  \cleardoublepage
  % 这部分的章节标题进行编号。
  \ctexset{
    chapter = {
      numbering = true,
    }
  }
  % 页码使用阿拉伯数字。
  \pagenumbering{arabic}
  \pagestyle{suepthesis}
  % 正文 22 磅的行距
  \setlength{\parskip}{0em}
  \linespread{1.53}\selectfont
  % 修复脚注出现跨页的问题
  \interfootnotelinepenalty=10000
}
%    \end{macrocode}
% \end{macro}
%
% \begin{macro}{\backmatter}
% 后置内容的页面样式。
%    \begin{macrocode}
\RenewDocumentCommand \backmatter {} {
  % 同样，所有的章节标题不进行编号。
  \setcounter{section}{0}
  \setcounter{subsection}{0}
  \setcounter{subsubsection}{0}
  \ctexset{
    chapter = {
      numbering = false,
      beforeskip = 18bp,
      format = {
        \heiti \l_@@_unnumchapter_style_cs:n \centering \zihao{3}
      },
      afterskip = 26bp,
    }
  }
}
%    \end{macrocode}
% \end{macro}
%
% 定义标题的前后间距。
%    \begin{macrocode}
\setlength{\abovecaptionskip}{11pt}
\@@_if_bachelor_thesis:TF {
  \@@_if_thesis_english:TF {
    \setlength{\belowcaptionskip}{9pt}
  } {
    % 为了满足 “前后一行空白的问题”，需要删除 caption 下方的间距。
    % 详见 `caption` 宏包手册和
    % https://github.com/CTeX-org/forum/issues/86
    %
    % 这里实际的 skip 在 15pt 左右，但是全部移除会导致当图片置于页面顶部时，
    % 图片与上方的间距过小，因此这里只移除 5pt。
    % 当然，这样会导致文本间的图片的 caption 下方的间距微微大于一行。
    %
    % 至于表格，虽然其caption位置在上方（而图片是在下方），
    % 但 `caption` 宏包已考虑这种区别，统一设置 `belowskip` 即可。
    \captionsetup{belowskip=-5pt}
    % 不过 longtable 宏包有另外的机制，不设置 belowskip 时间距已可较小，
    % 设置成负数还导致 caption 和表格本体的距离变大。因此我们撤销更改。
    \captionsetup[longtable]{belowskip=0pt}

    % 此外在浮动体内部，调整表格 caption 和表格本体间的距离。
    % 本来默认有一定空隙，现改为紧贴，这样更接近Word模板实作。
    \captionsetup[table]{skip=5pt}
  }
} {
  % 而研究生模板不存在这个问题。
  \setlength{\belowcaptionskip}{9pt}
}
%    \end{macrocode}
%
% 定义分隔字符。
%    \begin{macrocode}
\tl_set:Nn \g_@@_label_divide_char_tl {.}
% 研究生模板要求 "图序和图题间空1个中文字距"
\captionsetup[figure]{font=small,labelsep=quad}
% 研究生模板要求 "表序和题目间空1个中文字距"
\captionsetup[table]{font=small,labelsep=quad}
% 其它 caption 也参照上述格式
\captionsetup[lstlisting]{font=small,labelsep=quad}
\captionsetup[algorithm]{font=small,labelsep=quad}
%    \end{macrocode}
%
% \begin{macro}{\thefigure,\thetable,\theequation,\thelstlisting,\lstlistingname}
% 定义各种计数器的格式。
%    \begin{macrocode}
% 图片：五号字。
\cs_set:Npn \thefigure {\thechapter\g_@@_label_divide_char_tl\arabic{figure}}

% 表格：五号字。
\cs_set:Npn \thetable {\thechapter\g_@@_label_divide_char_tl\arabic{table}}

% equation
\cs_set:Npn \theequation {\thechapter\g_@@_label_divide_char_tl\arabic{equation}}

% code snippet
\AtBeginDocument{
  \cs_gset:Npn \thelstlisting {\thechapter\g_@@_label_divide_char_tl\arabic{lstlisting}}
  \cs_gset:Npn \lstlistingname {\c_@@_label_code_tl}

  % 算法变成「章节号-序号」
  % 为了减少修改，我们只适配按章编号的情况。
  % 针对 algorithm 宏包
  \@ifpackagewith{algorithm}{chapter}{
    \cs_gset:Npn \thealgorithm
    {\thechapter\g_@@_label_divide_char_tl\arabic{algorithm}}
  }{}
  % 针对 algorithm2e 宏包
  \@ifpackagewith{algorithm2e}{algochapter}{
    % 名字中的“cf”是指其作者 Christophe Fiorio。
    \cs_gset:Npn \thealgocf
    {\thechapter\g_@@_label_divide_char_tl\arabic{algocf}}
  }{}

  % 默认的情况下，保留公式和上下文的一定间距。（会比 Word 稍宽一些）
  \setlength{\abovedisplayskip}{\l_@@_style_math_above_display_skip_dim}
  \setlength{\abovedisplayshortskip}{\l_@@_style_math_above_display_skip_dim}
  \setlength{\belowdisplayskip}{\l_@@_style_math_below_display_skip_dim}
  \setlength{\belowdisplayshortskip}{\l_@@_style_math_below_display_skip_dim}
  % 调整浮动体与文字之间的距离
  \addtolength{\intextsep}{\l_@@_misc_float_separation_tl\baselineskip}
  \addtolength{\textfloatsep}{\l_@@_misc_float_separation_tl\baselineskip}
  % longtable 宏包有另外的机制，需专门调整
  \@ifpackageloaded{longtable}{
    \addtolength{\LTpre}{\l_@@_misc_float_separation_tl\baselineskip}
    \addtolength{\LTpost}{\l_@@_misc_float_separation_tl\baselineskip}
  }{}
  % 调整算法与文字之间的距离
  % 针对 algorithm2e 宏包
  \@ifpackageloaded{algorithm2e}{
    % 宏包手册介绍可自定义宏，再`\SetAlgoSkip`；我们为简洁，直接覆写。
    \renewcommand{\@algoskip}{\vspace{\l_@@_misc_algorithm_separation_tl}}
  }{}
}
%    \end{macrocode}
% \end{macro}
%
% 调整底层 TeX 排版引擎参数以保证所有段落能够很好地以两端对齐的方式呈现。
% 是的，这是祖传代码。
% 在英文模式下禁用，因为这段代码会禁用 hyphenation.
%    \begin{macrocode}
\bool_if:NF \l_@@_style_hyphen_bool {
  \hbadness=10000
  \tolerance=1
  \emergencystretch=\maxdimen
  \hyphenpenalty=10000
}
%    \end{macrocode}
%
% 自定义一个默认的 lstlisting 样式。
%    \begin{macrocode}
\definecolor{codegreen}{rgb}{0,0.6,0}
\definecolor{codegray}{rgb}{0.5,0.5,0.5}
\definecolor{codepurple}{rgb}{0.58,0,0.82}
\definecolor{backcolour}{rgb}{0.95,0.95,0.92}
\lstdefinestyle{examplestyle}{
    backgroundcolor=\color{backcolour},
    commentstyle=\color{codegreen},
    keywordstyle=\color{magenta},
    numberstyle=\tiny\color{codegray},
    stringstyle=\color{codepurple},
    basicstyle=\ttfamily\footnotesize,
    breakatwhitespace=false,
    breaklines=true,
    captionpos=b,
    keepspaces=true,
    numbers=left,
    numbersep=5pt,
    showspaces=false,
    showstringspaces=false,
    showtabs=false,
    tabsize=2
}
\lstset{style=examplestyle}
%    \end{macrocode}
%
% 调整插图目录与表格目录的标题。
%    \begin{macrocode}
\cs_set:Npn \listfigurename {\currentpdfbookmark{\c_@@_label_figure_tl}{ch:figures}\@@_get_const:N {figure}}
\cs_set:Npn \listtablename {\currentpdfbookmark{\c_@@_label_table_tl}{ch:tables}\@@_get_const:N {table}}
%    \end{macrocode}
%
% 预定义用户常用的证明环境。
%    \begin{macrocode}
\theoremstyle{plain}
  \newtheorem{algo}{\@@_get_const:N {algo}}[chapter]
  \newtheorem{them}{\@@_get_const:N {them}}[chapter]
  \newtheorem{lem}{\@@_get_const:N {lem}}[chapter]
  \newtheorem{prop}{\@@_get_const:N {prop}}[chapter]
  \newtheorem{cor}{\@@_get_const:N {cor}}[chapter]
  \newtheorem{axi}{\@@_get_const:N {axi}}[chapter]
\theoremstyle{definition}
  \newtheorem{defn}{\@@_get_const:N {defn}}[chapter]
  \newtheorem{conj}{\@@_get_const:N {conj}}[chapter]
  \newtheorem{exmp}{\@@_get_const:N {exmp}}[chapter]
  \newtheorem{case}{\@@_get_const:N {case}}
\theoremstyle{remark}
  \newtheorem{rem}{\@@_get_const:N {rem}}
  \renewcommand{\qedsymbol}{\ensuremath{\blacksquare}}
%    \end{macrocode}
%
% \begin{macro}{\@@_dunderline:nnn,\@@_dunderline:nn,\@@_dunderline:n}
% 用于渲染下划线。
%
% 参数如下：
% \begin{itemize}
%   \item \#1 位置，可选值为 \texttt{c}enter、\texttt{l}eft、\texttt{r}ight。
%   \item \#2 |dim| 长度。
%   \item \#3 |tl| 文字内容。
% \end{itemize}
%    \begin{macrocode}
\cs_new:Npn \@@_dunderline:nnn #1#2#3 {
  {\setbox0=\hbox{#3}\ooalign{\copy0\cr\rule[\dimexpr#1-#2\relax]{\wd0}{#2}}}
}
\cs_new:Npn \@@_dunderline:nn #1#2 {
  \@@_dunderline:nnn {#1} {1pt} {#2}
}
\cs_new:Npn \@@_dunderline:n #1 {
  \@@_dunderline:nnn {-10pt} {1pt} {#1}
}
% 遗留代码，等待重构。
\newcommand\dunderline[3][-1pt]{{%
  \setbox0=\hbox{#3}
  \ooalign{\copy0\cr\rule[\dimexpr#1-#2\relax]{\wd0}{#2}}}}
%    \end{macrocode}
% \end{macro}
%
% \begin{macro}{|@@_render_cover_entry:nn}
% 用于渲染封面的辅助函数。
%
% 参数如下：
% \begin{itemize}
%   \item \#1 |{token_list}| 为封面信息条目的名称。
%   \item \#2 |{token_list}| 为封面信息条目的内容。
% \end{itemize}
%
% 需要在 |\l_@@_cover_label_max_width_dim| 和 |\l_@@_cover_value_max_width_dim|
% 存储已经计算出来的最大宽度。
%    \begin{macrocode}
\cs_new:Npn \@@_render_cover_entry:nn #1#2 {
  \makebox[\l_@@_cover_label_max_width_dim][\l_@@_cover_label_align_tl]{
    \tl_if_blank:VTF #1 {} {#1\l_@@_cover_delimiter_tl}
  }
  \hspace{1ex}
  \@@_dunderline:nnn{\l_@@_cover_underline_offset_dim}
    {\l_@@_cover_underline_thickness_dim}{
    \makebox[\l_@@_cover_value_max_width_dim][\l_@@_cover_value_align_tl]{#2}
  }\par
}
%    \end{macrocode}
% \end{macro}
%
% \begin{macro}{|@@_get_text_width:Nn,\@@_get_text_width:NV}
% 计算 \#2 所占用的宽度，将结果存储在 \#1 中。
%
% 参数如下：
% \begin{itemize}
%   \item \#1 |dim| 存储宽度的变量。
%   \item \#2 |tl| 要计算宽度的文本。
% \end{itemize}
%    \begin{macrocode}
\cs_new:Npn \@@_get_text_width:Nn #1#2
  {
    \hbox_set:Nn \l_tmpa_box {#2}
    \dim_set:Nn #1 { \box_wd:N \l_tmpa_box }
  }
\cs_generate_variant:Nn \@@_get_text_width:Nn { NV }
%    \end{macrocode}
% \end{macro}
%
% \begin{macro}{\@@_get_max_text_width:NN}
% 从 \#2 中获取最大的文本宽度，然后设置到 \#1 中。
%
% 参数如下：
% \begin{itemize}
%   \item \#1: |dim| 用于存储最大宽度。
%   \item \#2: |seq| 用于存储文本。
% \end{itemize}
%    \begin{macrocode}
\cs_new:Npn \@@_get_max_text_width:NN #1#2
  {
% 这里用 |group| 确保局部变量不会被污染。
    \group_begin:
      \seq_set_eq:NN \l_@@_tmpa_seq #2
      \dim_zero_new:N \l_@@_tmpa_dim
      \bool_until_do:nn { \seq_if_empty_p:N \l_@@_tmpa_seq }
        {
          \seq_pop_left:NN \l_@@_tmpa_seq \l_@@_tmpa_tl
          \@@_get_text_width:NV \l_@@_tmpa_dim \l_@@_tmpa_tl
          % 在两边加上空白，避免文本太靠边。
          \dim_gset:Nn #1 { \dim_max:nn {#1} { \l_@@_tmpa_dim + \l_@@_cover_auto_width_padding_dim * 2} }
        }
    \group_end:
  }
%    \end{macrocode}
% \end{macro}
%
% \begin{macro}{\@@_parse_entry}
% 解析封面信息条目。
%
% 参数如下：
% \begin{itemize}
%   \item \#1: |tl| 为封面信息条目的名称。
%   \item \#2: |tl| 为封面信息条目的内容。
% \end{itemize}
% |\\| 会被视为换行符，从而实现信息条目换行的效果。
%
%    \begin{macrocode}
\cs_new:Npn \@@_parse_entry #1 #2 {
  \seq_set_split:Nnx \l_@@_tmp_right_seq {\\} {#2}
  \seq_clear:N \l_@@_tmp_left_seq
  \seq_map_inline:Nn \l_@@_tmp_right_seq {
    \seq_put_right:Nn \l_@@_tmp_left_seq {}
  }
  \seq_put_left:Nn \l_@@_tmp_left_seq {#1}
  \seq_pop_right:NN \l_@@_tmp_left_seq \g_@@_trashcan_tl
}
%    \end{macrocode}
% \end{macro}
%
% \begin{macro}{\@@_render_cover_entry}
% 渲染封面信息项。此函数为主函数。
%    \begin{macrocode}
\cs_new:Npn \@@_render_cover_entry:n #1 {
  % 左边是标签，右边是值。
  % 形如：
  % { {label_1} {value_1}, {label_2} {value 2} }
  % 首先转换成 seq 类型。
  \seq_set_from_clist:NN \l_@@_input_seq #1
  \seq_map_inline:Nn \l_@@_input_seq {
    % 然后对于每一对 label 和 value，首先查找
    % value 中是否含有 \\ 字符，如果有，则将其分割成多个
    % label - value 对。
    % 比如 {label_1} {value \\ 1} 会被转换成
    % { {label_1} {value}, {} {1} }
    \@@_parse_entry ##1
    % 然后将这些 label - value 对添加到 \l_@@_right_seq
    % 或者 \l_@@_left_sql 中。
    % left 就是 label，right 就是 value。
    \seq_concat:NNN \l_@@_right_seq \l_@@_right_seq \l_@@_tmp_right_seq
    \seq_concat:NNN \l_@@_left_seq \l_@@_left_seq \l_@@_tmp_left_seq
  }

  % 如果用户选择自动计算最大宽度，则计算最大宽度。
  \bool_if:NT \l_@@_cover_auto_width_bool {
    \@@_get_max_text_width:NN \l_@@_cover_label_max_width_dim \l_@@_left_seq
    \@@_get_max_text_width:NN \l_@@_cover_value_max_width_dim \l_@@_right_seq
  }


  % 最后，根据宽度渲染 label 和 value 对。
  \group_begin:
    \bool_until_do:nn { \seq_if_empty_p:N \l_@@_left_seq }
      {
        \seq_pop_left:NN \l_@@_left_seq \l_@@_tmpa_tl
        \seq_pop_left:NN \l_@@_right_seq \l_@@_tmpb_tl
        \tl_if_empty:xTF \l_@@_tmpb_tl {} {
          \@@_render_cover_entry:nn {\l_@@_tmpa_tl} {\l_@@_tmpb_tl}
        }
      }
  \group_end:
}
%    \end{macrocode}
% \end{macro}
%
% \begin{macro}{\@@_uline:n}
% 绘制下划线。
%    \begin{macrocode}
\cs_new:Npn \@@_uline:n #1
  { \CJKunderline [ thickness = 0.5pt ] { #1 } }
%    \end{macrocode}
% \end{macro}
%
% \begin{macro}{\@@_cover_iii:nnnn}
% \changes{v1.17.0.0}{2022/05/28}{绘制研究生中文题名页顶部信息}
% 绘制研究生中文题名页顶部信息。
% \begin{arguments}
%   \item 标签宽度。
%   \item 标签名称。
%   \item 值宽度。
%   \item 值内容。
% \end{arguments}
%    \begin{macrocode}
\cs_new:Npn \@@_cover_iii:nnnn #1#2#3#4
  {
    \dim_set:Nn \baselineskip { 20pt }
    \hbox:n
      {
        \rmfamily \zihao { 5 } \bfseries
        \hbox_to_wd:nn { #1 } { #2 }
        \skip_horizontal:n { 0.5em }
        \@@_uline:n { \hbox_to_wd:nn { #3 } { \hfil #4 \hfil } }
      }
  }
%    \end{macrocode}
% \end{macro}
%
% \begin{macro}{\make_graduate_cover:}
% 制作研究生论文模板封面。
%    \begin{macrocode}
\cs_new:Npn \make_graduate_cover: {
  \cleardoublepage
  \currentpdfbookmark{封面}{frontmatter:cover1}
  \begin{titlepage}
    { \songti\zihao{-4}
      \vbox:n { }
      \skip_vertical:n { -7.5pt }
      \dim_set:Nn \baselineskip { 15.6bp }
      \vbox:n
      {
        \mode_leave_vertical:
        \@@_cover_iii:nnnn { 4em } { 分~类~号 } { 7em } { \l_@@_clc_str }
        \hfill
        \@@_cover_iii:nnnn { 3em } { 密~级 } { 7em } { \l_@@_secret_lv_str }
      }
      \vbox:n
        {
          \mode_leave_vertical:
          \@@_cover_iii:nnnn { 4em } { 单位代码 } { 7em } { 10256 }
          \hfill
          \@@_cover_iii:nnnn { 3em } { 学~号 } { 7em } { \@@_secret_info:N{\l_@@_student_id_str}}
        }
    }

      \vspace{20mm}
      \begin{figure}
          \includegraphics[width=154.8pt,height=36.84pt]{images/graduate_logo.jpg}
      \end{figure}
      \vspace{20mm}

      \begin{center}
        \@@_if_academic_thesis:TF {
          \@@_if_master_thesis:TF {
            { 
              \heiti\zihao{-0}
              硕~士~学~位~论~文 \par
            } {
              \zihao{2}
              Dissertation~for~Master’s~Degree
            }
          } {
            { 
              \heiti\zihao{-0}
              博~士~学~位~论~文 \par
            } {
              \zihao{2}
              Dissertation~for~Doctor’s~Degree
            }
          }
  
        } {
          \@@_if_master_thesis:TF {
            { 
              \heiti\zihao{1}
              专业学位硕士学位论文 \par
            } {
              \zihao{-3}
              MASTER~THESIS~FOR~PROFESSIONAL~DEGREE
            }
          } {
            { 
              \heiti\zihao{1}
              专业学位博士学位论文\par
            } {
              \zihao{-3}
              DOCTOR~THESIS~FOR~PROFESSIONAL~DEGREE
            }
          }
        }
        \vspace*{65mm}
        {
          \heit\zihao{2}
          \l_@@_value_title_tl \par
          \l_@@_value_title_en_tl
        }
        \vspace*{65mm}
        {
          \renewcommand{\baselinestretch}{2}
  
          \tl_if_empty:NT \l_@@_cover_delimiter_tl {
            \tl_set:Nn \l_@@_cover_delimiter_tl {\qquad}
          }
          \tl_set:Nn \l_@@_cover_underline_offset_dim {-5pt}
  
          % 如果不是自动计算宽度，且用户没有自定义宽度，
          % 则尝试提供一个默认宽度。
          \bool_if:NF \l_@@_cover_auto_width_bool {
            \dim_compare:nNnT {\l_@@_cover_label_max_width_dim} = {0pt} {
              \dim_set:Nn \l_@@_cover_label_max_width_dim {45mm}
            }
            \dim_compare:nNnT {\l_@@_cover_value_max_width_dim} = {0pt} {
              \dim_set:Nn \l_@@_cover_value_max_width_dim {60mm}
            }
          }
  
          % 渲染信息。
          \clist_set:Nn \l_@@_input_clist {
              {\c_@@_graduate_label_author_tl} {\@@_secret_info:N{\l_@@_value_author_tl}},
              {\c_@@_graduate_label_supervisor_tl} {\@@_secret_info:N{\l_@@_value_supervisor_tl}},
              {\g_@@_graduate_label_major_tl} {\l_@@_value_major_tl},
              {\c_@@_graduate_label_degreeclass_tl} {\@@_secret_info:N{\l_@@_value_degree_class_tl}},
              {\c_@@_graduate_label_institute_tl} {\l_@@_value_institute_tl},
           }
  
          \heiti\zihao{-3}
          \@@_render_cover_entry:n \l_@@_input_clist
        }
        {\songti\zihao{4} \l_@@_cover_date_tl}
      \end{center}
  \end{titlepage}
}
%    \end{macrocode}
% \end{macro}
%
% \begin{macro}{\@@_make_gradute_title_page:}
% 制作中文封面页。
%    \begin{macrocode}
\cs_new:Npn \@@_make_gradute_title_page: {
  \cleardoublepage
  \currentpdfbookmark{中文扉页}{frontmatter:titlepage}
  \begin{titlepage}
    { \songti\zihao{-4}
      \vbox:n { }
      \skip_vertical:n { -7.5pt }
      \dim_set:Nn \baselineskip { 15.6bp }
      \vbox:n
      {
        \mode_leave_vertical:
        \@@_cover_iii:nnnn { 4em } { 分\quad类\quad号 } { 7em } { \l_@@_clc_str }
        \hfill
        \@@_cover_iii:nnnn { 3em } { 密\quad级 } { 7em } { \l_@@_secret_lv_str }
      }
      \vbox:n
        {
          \mode_leave_vertical:
          \@@_cover_iii:nnnn { 4em } { 单位代码 } { 7em } { 10256 }
          \hfill
          \@@_cover_iii:nnnn { 3em } { 学\quad号 } { 7em } { \@@_secret_info:N{\l_@@_student_id_str}}
        }
    }
      \vspace{20mm}

      \begin{center}
        \@@_if_academic_thesis:TF {
          \heiti\zihao{-1}
          \@@_if_master_thesis:TF {
            上海电力大学硕士学位论文 \par
          } {
            上海电力大学博士学位论文 \par
          }
  
        } {
          \heiti\zihao{-1}
          \@@_if_master_thesis:TF {
            上海电力大学 \\
            专业学位硕士学位论文 \par
            
          } {
            上海电力大学 \\
            专业学位士学位论文 \par
          }
        }
        \vspace*{65mm}
        {
          \heit\zihao{2}
          \l_@@_value_title_tl \par
          \l_@@_value_title_en_tl
        }
        \vspace*{65mm}
        {
          \renewcommand{\baselinestretch}{2}
  
          \tl_if_empty:NT \l_@@_cover_delimiter_tl {
            \tl_set:Nn \l_@@_cover_delimiter_tl {\qquad}
          }
          \tl_set:Nn \l_@@_cover_underline_offset_dim {-5pt}
  
          % 如果不是自动计算宽度，且用户没有自定义宽度，
          % 则尝试提供一个默认宽度。
          \bool_if:NF \l_@@_cover_auto_width_bool {
            \dim_compare:nNnT {\l_@@_cover_label_max_width_dim} = {0pt} {
              \dim_set:Nn \l_@@_cover_label_max_width_dim {45mm}
            }
            \dim_compare:nNnT {\l_@@_cover_value_max_width_dim} = {0pt} {
              \dim_set:Nn \l_@@_cover_value_max_width_dim {60mm}
            }
          }
  
          % 渲染信息。
          \clist_set:Nn \l_@@_input_clist {
              {\c_@@_graduate_label_author_tl} {\@@_secret_info:N{\l_@@_value_author_tl}},
              {\c_@@_graduate_label_supervisor_tl} {\@@_secret_info:N{\l_@@_value_supervisor_tl}},
              {\g_@@_graduate_label_major_tl} {\l_@@_value_major_tl},
              {\c_@@_graduate_label_degreeclass_tl} {\@@_secret_info:N{\l_@@_value_degree_class_tl}},
              {\c_@@_graduate_label_finalizationdate_tl} {\l_@@_value_finalization_date_tl},
           }
          \heiti\zihao{-3}
          \@@_render_cover_entry:n \l_@@_input_clist
        }
      \end{center}
  \end{titlepage}
}
%    \end{macrocode}
% \end{macro}
%
% \begin{macro}{\circled}
% 圆形数字编号定义。
%    \begin{macrocode}
\newcommand{\circled}[2][]{\tikz[baseline=(char.base)]
  {\node[shape = circle, draw, inner~sep = 1pt]
  (char) {\phantom{\ifblank{#1}{#2}{#1}}};
  \node at (char.center) {\makebox[0pt][c]{#2}};}}
\robustify{\circled}
%    \end{macrocode}
% \end{macro}
%
% \begin{macro}{\@@_graduate_originality:}
% 研究生原创性声明。
%    \begin{macrocode}
\cs_new:Npn \@@_graduate_originality:
  {
    % 取消页眉页脚。
    \ctexset {
      chapter / pagestyle = plain,
    }

    \begin{titlepage}
      % 不计算页码。
      \pagenumbering{gobble}

      % 原创性声明部分
      \begin{center}
        \@@_same_page:
        \ctexset{
          chapter = {
            titleformat = {\heiti\zihao{-2}},
          }
        }
        \currentpdfbookmark{\c_@@_graduate_label_originality_tl}{frontmatter:originality}
        \chapter*{
          \c_@@_graduate_label_originality_tl
        }
      \end{center}

      % 本部分字号为三号。
      \zihao{3}
      \qquad\c_@@_graduate_label_originality_clause_tl
      \vspace{3\baselineskip}

      \begin{flushright}
        \c_@@_graduate_label_originality_author_signature_tl\par
      \end{flushright}

      \vspace{3\baselineskip}

      % 使用授权声明部分。
      \begin{center}
        \@@_same_page:
        \ctexset{
          chapter = {
            titleformat = {\heiti\zihao{-2}},
          }
        }
        \currentpdfbookmark{\c_@@_graduate_label_authorization_tl}{frontmatter:originality1}
        \chapter*{
          \c_@@_graduate_label_authorization_tl
        }
      \end{center}

      \qquad\c_@@_graduate_label_authorization_clause_tl

      \vspace*{15mm}

      \begin{flushright}
        \begin{spacing}{1.65}
          \zihao{3}
          % \hspace{5mm}\raisebox{-2ex}{\includegraphics[width=30mm]{example-image}}\hspace{5mm}
          \c_@@_graduate_label_originality_author_signature_tl\par
          \vspace{1\baselineskip}
          \c_@@_graduate_label_originality_supervisor_signature_tl\par
        \end{spacing}
      \end{flushright}
    \end{titlepage}
    \cleardoublepage
  }
%    \end{macrocode}
% \end{macro}
%
% \subsubsection{定义用户接口}
%
% \begin{macro}{\SUEPSetup}
% 提供用户配置的接口。
%    \begin{macrocode}
\DeclareDocumentCommand \SUEPSetup { m }
  { \keys_set:nn { suepthesis } { #1 }}
%    \end{macrocode}
% \end{macro}
%
% \begin{macro}{\BigStar}
% 提供密级选项中需要的五角星，在普通环境中使用。
%    \begin{macrocode}
\DeclareDocumentCommand \BigStar { }
  { \ding{72} }
%    \end{macrocode}
% \end{macro}
%
% \begin{environment}{blindPeerReview}
% 用于包裹涉及个人信息的内容。
%
% 在启用盲审模式时，其中的内容会被隐藏。
%
% 本环境提供了一个可选参数，可以传入一个 bool 值，用于在盲审模式下关闭
% 隐藏行为。
%    \begin{macrocode}
  \NewDocumentEnvironment {blindPeerReview} {O{\c_true_bool} +b}
  {
    \bool_if:nTF {\g_@@_blind_mode_bool && #1} {} {
      #2
    }
  } {}
%    \end{macrocode}
% \end{environment}
%
% \begin{macro}{\cleardoublepage}
% 重定义 \tn{cleardoublepage}，
% 使得偶数页面在没有内容时也不显示页眉页脚。见：
% \url{https://tex.stackexchange.com/a/1683}。
%    \begin{macrocode}
\RenewDocumentCommand \cleardoublepage { }
  {
    \clearpage
    \bool_if:NT \g_@@_twoside_bool
      {
        \int_if_odd:nF \c@page
          { \hbox:n { } \thispagestyle { empty } \newpage }
      }
  }
%    \end{macrocode}
% \end{macro}
%
% \begin{macro}{\SecretInfo}
% 用于包裹涉及个人信息的内容。
%   \begin{macrocode}
\DeclareDocumentCommand \SecretInfo { m o }
  {
    \IfValueTF {#2} {
      \@@_secret_info:nn {#1} {#2}
    } {
      \@@_secret_info:n {#1}
    }
  }
%   \end{macrocode}
% \end{macro}
%
% \begin{macro}{\MakeCover}
% 制作封面。
%    \begin{macrocode}
\DeclareDocumentCommand \MakeCover {}
  {
    
    \begin{blindPeerReview}[\l_@@_cover_hide_cover_in_peer_review_bool]
      \@@_if_graduate:TF {
        \make_graduate_cover:
        \@@_make_gradute_title_page:
      } {

      }
    \end{blindPeerReview}
  }
%    \end{macrocode}
% \end{macro}
%
%
% \begin{macro}{\MakeOriginality}
% 原创性声明。
%    \begin{macrocode}
\NewDocumentCommand \MakeOriginality {}
  {
    \group_begin:
    \begin{blindPeerReview}[\l_@@_cover_hide_cover_in_peer_review_bool]
      \int_case:nn {\g_@@_thesis_type_int}
      {
        {1}
        {
          \currentpdfbookmark{声明}{frontmatter:originality}
          \pagestyle{suepthesis}
          \pagenumbering{gobble}

          % 原创性声明部分
          \begin{center}
            \vspace*{-2bp}
            \@@_same_page:
            \chapter*{\heiti\zihao{2}\c_@@_bachelor_label_originality_tl}
          \end{center}~\par

          % 本部分字号为小三。
          \zihao{-3}
          \c_@@_bachelor_label_originality_clause_tl

          \vspace{17mm}

          \begin{flushright}
            \c_@@_bachelor_label_originality_author_signature_tl\par
          \end{flushright}

          \vspace{16mm}

          % 使用授权声明部分
          \begin{center}
            \@@_same_page:
            \chapter*{
              \heiti\zihao{2}
              \c_@@_bachelor_label_authorization_tl
            }
          \end{center}~\par

          \c_@@_bachelor_label_authorization_clause_tl

          \vspace*{3mm}

          \begin{flushright}
            \begin{spacing}{1.65}
              \zihao{-3}
              % \hspace{5mm}\raisebox{-2ex}{\includegraphics[width=30mm]{example-image}}\hspace{5mm}
              \c_@@_bachelor_label_originality_author_signature_tl\par
              \c_@@_bachelor_label_originality_supervisor_signature_tl\par
            \end{spacing}
          \end{flushright}

          \newpage
        }
        {3} {
          \currentpdfbookmark{Statements}{frontmatter:originality}
          \pagestyle{suepthesis}
          \pagenumbering{gobble}

          % 原创性声明部分
          \begin{center}
            \vspace*{-2bp}
            \@@_same_page:
            \chapter*{
              \heiti\zihao{-2}
              \c_@@_bachelor_english_label_originality_tl
            }
          \end{center}~\par

          % 本部分字号为小四
          \zihao{-4}
          \c_@@_bachelor_english_label_originality_clause_tl

          \bigbreak

          Student~(Signature):~\dunderline[-1pt]{1pt}{\makebox[18mm]{}}~Date:\par

          \vspace{\stretch{1}}

          % 使用授权声明部分
          \begin{center}
            \@@_same_page:
            \chapter*{
              \heiti\zihao{-2}
              \c_@@_bachelor_english_label_authorization_tl
            }
          \end{center}~\par

          \c_@@_bachelor_english_label_authorization_clause_tl

          \bigbreak
          Student~(Signature):~
            \dunderline[-1pt]{1pt}{\makebox[18mm + 16bp]{}}~
            \hspace{2mm}Date:\par
          Supervisor~(Signature):~
            \dunderline[-1pt]{1pt}{\makebox[18mm]{}}~
            \hspace{2mm}Date:\par
        }
        {4} {\@@_graduate_originality:}
        {5} {\@@_graduate_originality:}
      }
    % 单独成页
    \clearpage
    \end{blindPeerReview}
    \group_end:
  }
%    \end{macrocode}
% \end{macro}
%
% \begin{macro}{\MakeTOC}
% 生成目录。
%    \begin{macrocode}
\DeclareDocumentCommand \MakeTOC {}
  {
    {
      \@@_if_bachelor_thesis:TF {
        \renewcommand{\baselinestretch}{1.35}
      } {
        \renewcommand{\baselinestretch}{1.56}
      }

      % 自定义目录样式
      \cs_set:Npn \contentsname {
        \fontsize{16pt}{\baselineskip}
        \l_@@_unnumchapter_style_cs:n
          \heiti
            {\l_@@_toc_title_tl}
        \vspace{-8pt}
      }

      \bool_if:NTF \l_@@_add_toc_to_toc_bool {
        % 添加「目录」本身到目录中，同时自动添加书签
        % 此处必须有`\phantomsection`，不然 hyperref 会把链接指向之前摘要的标题。
        \phantomsection
        \addcontentsline{toc}{chapter}{\c_@@_label_toc_en_tl}
      } {
        % 手动添加目录书签
        \currentpdfbookmark{\l_@@_toc_title_tl}{ch:toc}
      }

      % 制作目录
      \tableofcontents

      % 单独成页
      \clearpage
    }
  }
%    \end{macrocode}
% \end{macro}
%
% \begin{environment}{abstract}
% 生成摘要。
%    \begin{macrocode}
\NewDocumentEnvironment {abstract} {}
  {

    \cleardoublepage
    \linespread{1.53}\selectfont

    \@@_if_bachelor_thesis:T {
      \begin{center}
        \vspace*{-17bp}
        \heiti\zihao{-2}\textbf{
          \int_case:nn {\g_@@_thesis_type_int}
          {
            {1} {\l_@@_value_title_tl}
            {2} {\l_@@_value_trans_title_tl}
            {3} {\l_@@_value_title_tl}
          }
        }
      \end{center}

      \vspace*{2mm}
    }

    \ctexset{
      chapter/numbering = false,
    }

    \@@_if_bachelor_thesis:T {
      \ctexset{
        chapter/titleformat = {\textmd}
      }
    }

    {
      \@@_same_page:
      \bool_if:NTF \l_@@_add_abstract_to_toc_bool {
        \chapter{\c_@@_label_abstract_tl}
      } {
        \currentpdfbookmark{\c_@@_label_abstract_tl}{ch:abstract}
        \chapter*{\c_@@_label_abstract_tl}
      }
    }
    \vspace*{1mm}
    \par
  }
  {
    \par
    \vspace{4ex}
    \noindent
    \@@_if_graduate:TF {
      % 研究生模板中，“关键词”宋体小四加粗
      % 关键词为宋体小四号字。
      \textbf{\c_@@_label_keywords_tl}\l_@@_value_keywords_tl\par
    } {
      % 本科生模板中，关键词为黑体加粗
      \textbf{\heiti \c_@@_label_keywords_tl \l_@@_value_keywords_tl}\par
    }
    \newpage
  }
%    \end{macrocode}
% \end{environment}
%
% \begin{environment}{abstractEn}
% 生成英文摘要。
%    \begin{macrocode}
\NewDocumentEnvironment {abstractEn} {}
  {
    \linespread{1.53}\selectfont

    \@@_if_bachelor_thesis:T {
      \begin{spacing}{0.95}
        \centering
        \vspace*{-2bp}

        \heiti {
          \zihao{3}\textbf
          \l_@@_value_title_en_tl\\
        }
      \end{spacing}
      \vspace*{10mm}
    }

    \ctexset{
      chapter/numbering = false,
    }

    \@@_if_bachelor_thesis:TF {
      \ctexset{
        chapter/titleformat = {\zihao{-3}\textmd}
      }
    } {
      \ctexset {
        chapter/titleformat = {\heiti\zihao{3}\centering\textbf}
      }
    }

    {
      \@@_same_page:
      \bool_if:nTF {\l_@@_add_abstract_en_to_toc_bool} {
        \chapter{\c_@@_label_abstract_en_tl}
      } {
        \currentpdfbookmark{\c_@@_label_abstract_en_tl}{ch:abstract:en}
        \chapter*{\c_@@_label_abstract_en_tl}
      }
    }
  }
  {
    \par\vspace{3ex}\noindent
    \@@_if_graduate:TF {
      % Times New Roman小四号字，行距22磅
      % “Key Words”
      % Times New Roman小四号字加粗
      \textbf{\c_@@_label_keywords_en_tl} \l_@@_value_keywords_en_tl
    } {
      \textbf{\c_@@_label_keywords_en_tl \l_@@_value_keywords_en_tl}
    }
    \newpage
  }

%    \end{macrocode}
% \end{environment}
%
% \begin{environment}{conclusion}
% 生成结论。需要放在 \cs{macrocode} 之后。
%    \begin{macrocode}
\NewDocumentEnvironment {conclusion} {}
  {
    \ctexset{
      section/number = \arabic{section}
    }

    \@@_if_thesis_english:TF {
      \chapter{\c_@@_label_conclusion_en_tl}
    } {
      \chapter{\c_@@_label_conclusion_tl}
    }
  }
  {}
%    \end{macrocode}
% \end{environment}
%
% \begin{environment}{bibprint}
% 生成参考文献。需要放在 \cs{backmatter} 之后。
%    \begin{macrocode}
\NewDocumentEnvironment {bibprint} {}
  {
    % 设置参考文献字号为 5 号
    \renewcommand*{\bibfont}{\zihao{5}}
    % 设置参考文献各个项目之间的垂直距离为 0
    \setlength{\bibitemsep}{0ex}
    \setlength{\bibnamesep}{0ex}
    \setlength{\bibinitsep}{0ex}
    \@@_if_graduate:TF {
    } {
      % 「本科生」设置单倍行距
      \renewcommand{\baselinestretch}{1.2}
    }
    % 设置参考文献顺序标签 `[1]` 与文献内容 `作者. 文献标题...` 的间距
    \setlength{\biblabelsep}{1.7mm}

    \bool_if:NF \l_@@_style_bibliography_indent_bool {
      % 设置参考文献后文缩进为 0（与 Word 模板保持一致）
      % See: https://github.com/hushidong/biblatex-gb7714-2015
      %      如何修参考文献表的缩进？
      \cs_set:Npn \itemcmd {
        \settowidth{\lengthid}{\mkgbnumlabel{\printfield{labelnumber}}}
        %%这里是所做的调整，以下两句通过调整\lengthid来调整缩进
        \setlength{\lengthid}{0pt}
        \addtolength{\lengthid}{-\biblabelsep}
        \setlength{\lengthlw}{\textwidth}
        \addtolength{\lengthlw}{-\lengthid}
        \addvspace{\bibitemsep}%恢复\bibitemsep的作用
        \hangindent\lengthid
        \leavevmode\mkgbnumlabel{\printfield{labelnumber}}%
        \hspace{\biblabelsep}
      }
    }

    \@@_if_thesis_english:TF {
      \chapter{\c_@@_label_reference_en_tl}
    } {
      \chapter{\c_@@_label_reference_tl}
    }
  }
  {}
%    \end{macrocode}
% \end{environment}
%
% \begin{environment}{appendices}
% 生成附录。
%    \begin{macrocode}
\NewDocumentEnvironment {appendices} {}
  {
    % Used in chapter, ToC.
    \tl_new:N \l_@@_appendix_plain_label_tl
    % Used before reference label.
    \tl_new:N \l_@@_appendix_default_title_tl

    \@@_if_thesis_english:TF {
      \tl_set:Nn \l_@@_appendix_plain_label_tl {\c_@@_label_appendix_prefix_en_tl}
      \tl_set:Nn \l_@@_appendix_default_title_tl {\c_@@_label_appendix_en_tl}
    } {
      \tl_set:Nn \l_@@_appendix_plain_label_tl {\c_@@_label_appendix_prefix_tl}
      \tl_set:Nn \l_@@_appendix_default_title_tl {\c_@@_label_appendix_tl}
    }

    \bool_if:NTF \l_@@_appendices_chapter_level_bool {
      % 使用以「chapter」为顶层的附录格式

      % 仅设置 \setcounter{chapter}{0} 时，pdf 目录会索引到正文章节。
      % 因此，需要使用 \appendix 重置计数器，并将附录后面的
      % 几个章节视为特殊的附录页。
      \appendix

      \ctexset{
        chapter/numbering = true,
        chapter/name = {},
        chapter/number = \l_@@_appendix_plain_label_tl\hspace{1ex}\Alph{chapter},
        section/number = \Alph{chapter}. \arabic{section},
        subsection/number = \Alph{chapter}. \arabic{section}. \arabic{subsection},
      }

      \cs_set:Npn \thechapter {
        \Alph{chapter}
      }
    } {
      % 使用以「section」为顶层的附录格式

      % 因为不需要用到 chapter counter, 所以直接加一即可。
      \stepcounter{chapter}
      \setcounter{section}{0}
      % （与上面方法至少用一个）
      % 需要让 section 在 pdf bookmark 中输出字母而不是数字。
      % 详见 hyperref 代码。
      \gdef\theHsection{\Alph{section}}

      % 定义 \thefigure 采用节而不是章
\cs_set:Npn \thefigure {\theHsection \g_@@_label_divide_char_tl\arabic{figure}}

      \ctexset{
        section/number = \l_@@_appendix_plain_label_tl\hspace{1ex}\Alph{section},
        subsection/number = \Alph{section}. \arabic{subsection},
      }

      \cs_gset:Npn \thechapter {
        \Alph{section}
      }
      % \gdef \thechapter{\Alph{section}}

      \tl_if_blank:VTF \l_@@_appendices_title_tl {
        \chapter{\l_@@_appendix_default_title_tl}
      } {
        \chapter*{\l_@@_appendices_title_tl}
        \stepcounter{chapter}
        \tl_if_blank:VTF \l_@@_appendix_toc_title_tl {
          \addcontentsline{toc}{chapter}{\l_@@_appendix_default_title_tl}
        } {
          \addcontentsline{toc}{chapter}{\l_@@_appendix_toc_title_tl}
        }
      }
    }
  }
  {
  }
%    \end{macrocode}
% \end{environment}
%
% \begin{environment}{acknowledgements}
% 生成致谢。
%    \begin{macrocode}
\NewDocumentEnvironment {acknowledgements} {+b}
  {
    \begin{blindPeerReview}
      % 将此章节视为特殊的附录页，关闭附录编号，重定义 section 编号。
      % 不知为何，需要手动重置 section 计数器。
      \setcounter{section}{0}
      \ctexset{
        appendix/numbering = false,
        section/number = \arabic{section},
        subsection/number = \arabic{section}. \arabic{subsection},
        subsubsection/number = \arabic{section}. \arabic{subsection}. \arabic{subsubsection},
      }

      \chapter{\g_@@_const_heading_acknowledgements_tl}
      \@@_if_graduate:TF {\fangsong}{}
      #1
    \end{blindPeerReview}
  } {}
%    \end{macrocode}
% \end{environment}
%
% \begin{macro}{\Author,\AuthorEn}
% 在普通模式下，输出作者姓名。
% 在盲审模式下，输出「第 n 作者」。
%    \begin{macrocode}
\NewDocumentCommand \Author {O{1} o o}
  {
    \bool_if:NTF \g_@@_blind_mode_bool {
      % 盲审模式
      \IfValueTF {#3} {
        #3
      } {
        第\zhnumber{#1}作者
      }
    } {
      % 普通模式
      \IfValueTF {#2} {
        % 覆盖默认的 \author 命令
        #2
      } {
        % 默认采用作者姓名
        \l_@@_value_author_tl
      }
    }
  }

% 英文姓名
\NewDocumentCommand \AuthorEn {O{1} o o}
  {
    \bool_if:NTF \g_@@_blind_mode_bool {
      % 盲审模式
      \IfValueTF {#3} {
        #3
      } {
        \Ordinalstringnum{#1}~Author
      }
    } {
      % 普通模式
      \IfValueTF {#2} {
        % 覆盖默认的 \author 命令
        #2
      } {
        % 默认采用作者姓名
        \l_@@_value_author_en_tl
      }
    }
  }
%    \end{macrocode}
% \end{macro}
%
% \begin{macro}{\addpub,\addpubs}
% 添加一个或多个参考文献。
%   \begin{macrocode}
\NewDocumentCommand \addpub {m} {
  \nocite{#1}
  \addtocategory{mypub}{#1}
}

\NewDocumentCommand \addpubs {m} {
  % apply a clist
  \clist_map_function:nN {#1} \addpub
}
%   \end{macrocode}
% \end{macro}
%
% \begin{macro}{\pubsection}
% 设置小标题。
%    \begin{macrocode}
\NewDocumentCommand \pubsection {m} {
  {
    % 自增计数器
    \par\stepcounter{pub}
    % 设置小标题，暂时没有考虑英文模式
    \noindent
    \textbf{
      \heiti{
        \zhnumber{\thepub}、#1
      }
    }\par
  }
}
%    \end{macrocode}
% \end{macro}
%
% \begin{environment}{publications}
% 生成攻读学位期间发表论文与研究成果清单。
%    \begin{macrocode}
\NewDocumentEnvironment {publications} {+b}
  {
    % 同时设置 omit 以及 blindPeerReview 才能跳过此章节生成。
    \begin{blindPeerReview}[\l_@@_publications_omit_bool]
      % 将此章节视为特殊的附录页，关闭附录编号，重定义 section 编号。
      % 不知为何，需要手动重置 section 计数器。
      \setcounter{section}{0}
      \ctexset{
        appendix/numbering = false,
        section/number = \arabic{section},
        subsection/number = \arabic{section}. \arabic{subsection},
        subsubsection/number = \arabic{section}. \arabic{subsection}. \arabic{subsubsection},
      }
      % 设置参考文献字号为 5 号
      \renewcommand*{\bibfont}{\zihao{5}}
      % 设置参考文献各个项目之间的垂直距离为 0
      \setlength{\bibitemsep}{0ex}
      \setlength{\bibnamesep}{0ex}
      \setlength{\bibinitsep}{0ex}
      % 设置参考文献顺序标签 `[1]` 与文献内容 `作者. 文献标题...` 的间距
      \setlength{\biblabelsep}{1.7mm}

      \bool_if:NF \l_@@_style_bibliography_indent_bool {
        % 设置参考文献后文缩进为 0（与 Word 模板保持一致）
        % See: https://github.com/hushidong/biblatex-gb7714-2015
        %      如何修参考文献表的缩进？
        \cs_set:Npn \itemcmd {
          \settowidth{\lengthid}{\mkgbnumlabel{\printfield{labelnumber}}}
          %%这里是所做的调整，以下两句通过调整\lengthid来调整缩进
          \setlength{\lengthid}{0pt}
          \addtolength{\lengthid}{-\biblabelsep}
          \setlength{\lengthlw}{\textwidth}
          \addtolength{\lengthlw}{-\lengthid}
          \addvspace{\bibitemsep}%恢复\bibitemsep的作用
          \hangindent\lengthid
          \leavevmode\mkgbnumlabel{\printfield{labelnumber}}%
          \hspace{\biblabelsep}
        }
      }

      % If in blindPeerReview mode, omit delimiters in author field.
      \bool_if:NT \g_@@_blind_mode_bool {
        % 如果有多个作者，不修改此项的话，作者与标题之间会有逗号。
        \DeclareDelimFormat[bib,biblist]{finalnamedelim}{}
        % 如果自己不是第一个作者，不修改此项的话，会在最开始有逗号。
        \DeclareDelimFormat{multinamedelim}{}
        % 如果覆盖的是英文作者，不修改此项的话，会在最开始有空格。
        \DeclareDelimFormat{bibnamedelimd}{}

        % 如果作者太多而被截断，不修改的话，会多余逗号、“等”云云。
        % 被截断的充要条件：作者数量大于通过`\BITSetup`设置的`publications/maxbibnames`。
        % 故设置标点来去掉逗号，
        \DeclareDelimFormat[bib,biblist]{andothersdelim}{}
        % 并设置本地化字符串来去掉“等”。
        \setlocalbibstring{andothers}{}
        \setlocalbibstring{andotherscn}{}
        % 另外注意，我们仍尊重`maxbibnames`和`minbibnames`，保证若开盲审时显示作者，则关盲审时也正常。
      }

      % ===== 上方定义与「参考文献」部分相同

      % 中文姓名下，此部分不参与输出。
      \cs_set:Npn \mkbibnamegiven ##1 {
        \haspartannotation{myself}{
          \bool_if:NTF \g_@@_blind_mode_bool {
            % 盲审模式，不输出内容
          } {
            % 普通模式
            \textbf{##1}
          }
        }{
          \bool_if:NTF \g_@@_blind_mode_bool {
            % 盲审模式，不输出内容
          } {
            % 普通模式
            ##1
          }
        }
      }

      \cs_set:Npn \mkbibnamefamily ##1 {
        \haspartannotation{myself}{
          % 作者为自己
          \bool_if:NTF \g_@@_blind_mode_bool {
            % 盲审模式
            \getpartannotation{myself}
          } {
            % 普通模式
            \textbf{##1}
          }
        }{
          % 作者不是自己
          \bool_if:NTF \g_@@_blind_mode_bool {
            % 盲审模式，不输出
          } {
            % 普通模式
            ##1
          }
        }
      }

      \if_cs_exist:N \c@pub {
        % 重置计数器
        \setcounter{pub}{0}
      } \else: {
        % 设置计数器
        \newcounter{pub}
      } \fi:

      % 设置参考文献的排序
      \bool_if:NTF \l_@@_publications_sorting_bool {
        % Sorting by year, name, type.
        \newrefcontext[sorting=ynt]
      } {
        % Do not sort.
        \newrefcontext
      }

      % 根据 maxbibnames 的设置，覆盖 \blx@maxbibnames 选项，保证所有作者都能显示。
      \cs_set:Npn \blx@maxbibnames {
        \l_@@_publications_maxbibnames_int
      }

      % 根据 minbibnames 的设置，覆盖 \blx@minbibnames 选项，保证所有作者都能显示。
      \cs_set:Npn \blx@minbibnames {
        \l_@@_publications_minbibnames_int
      }

      \chapter{\@@_get_const:N {publications}}
      #1
    \end{blindPeerReview}
  }
  {}
%    \end{macrocode}
% \end{environment}
%